
\documentclass[11pt,a4paper]{article}
\usepackage[utf8]{inputenc}
\usepackage[T1]{fontenc}
\usepackage[french]{babel}
\usepackage{lmodern}
\usepackage{microtype}
\usepackage{amsmath,amssymb,amsthm,mathtools}
\usepackage{geometry}
\usepackage{fancyhdr}
\usepackage{lastpage}
\usepackage{xcolor}
\usepackage[most]{tcolorbox}
\usepackage{enumitem}
\usepackage{hyperref}
\usepackage{multicol}
\usepackage{array,booktabs}
\usepackage{needspace}
\geometry{margin=2.05cm,headheight=14pt}
\definecolor{bleuprepa}{RGB}{20,55,110}
\definecolor{vertprepa}{RGB}{20,105,70}
\definecolor{rougeprepa}{RGB}{160,35,35}
\definecolor{orangeprepa}{RGB}{180,100,20}
\definecolor{grisclair}{RGB}{245,247,250}
\hypersetup{colorlinks=true, linkcolor=bleuprepa, urlcolor=bleuprepa, citecolor=bleuprepa}
\pagestyle{fancy}
\fancyhf{}
\lhead{MPSI--PCSI--MP}
\chead{TD Analyse 4 -- Dérivation}
\rhead{Pierre Garcia}
\cfoot{\thepage/\pageref{LastPage}}
\setlength{\parindent}{0pt}
\setlength{\parskip}{0.45em}
\setlist[enumerate]{label=\textbf{\arabic*.}, leftmargin=1.1cm, itemsep=0.35em, topsep=0.35em}
\setlist[itemize]{leftmargin=0.8cm, itemsep=0.25em, topsep=0.25em}
\tcbset{enhanced,breakable,sharp corners,boxrule=0.7pt,colback=white,fonttitle=\bfseries}
\newtcolorbox{objectifbox}{colframe=bleuprepa,colback=blue!3!white,title=Objectif}
\newtcolorbox{methodebox}{colframe=vertprepa,colback=green!3!white,title=Méthode}
\newtcolorbox{attentionbox}{colframe=rougeprepa,colback=red!3!white,title=Attention}
\newtcolorbox{exobox}[2]{colframe=bleuprepa,colback=white,title={Exercice #1 -- #2}}
\newtcolorbox{corrbox}[2]{colframe=vertprepa,colback=white,title={Correction #1 -- #2}}
\newcommand{\R}{\mathbb{R}}
\newcommand{\N}{\mathbb{N}}
\newcommand{\Z}{\mathbb{Z}}
\newcommand{\dd}{\mathrm{d}}
\newcommand{\e}{\mathrm{e}}
\newcommand{\ch}{\operatorname{ch}}
\newcommand{\sh}{\operatorname{sh}}
\newcommand{\thh}{\operatorname{th}}
\newcommand{\Arctan}{\operatorname{arctan}}
\newcommand{\Arcsin}{\operatorname{arcsin}}
\newcommand{\Arccos}{\operatorname{arccos}}
\newcommand{\etoiles}[1]{\textbf{Difficulté :} #1}
\newcommand{\intervalle}[2]{\left[#1,#2\right]}
\begin{document}
\begin{titlepage}
\centering
\vspace*{1cm}
{\Large\bfseries Classe préparatoire scientifique}\par
\vspace{0.4cm}
{\Huge\bfseries TD complet d'analyse}\par
\vspace{0.25cm}
{\LARGE\bfseries Chapitre 4 -- Dérivation, théorèmes fondamentaux et fonctions usuelles}\par
\vspace{0.7cm}
{\large Niveau MPSI / PCSI / MP}\par
\vfill
\begin{tcolorbox}[colframe=bleuprepa,colback=blue!3!white,width=0.86\textwidth,title=Programme travaillé]
Dérivabilité en un point.
Dérivées latérales.
Approximation affine locale.
Tangente.
Opérations sur les dérivées.
Dérivée d'une composée.
Dérivée d'une réciproque.
Extrema locaux.
Théorème de Rolle.
Théorème des accroissements finis.
Inégalité des accroissements finis.
Logarithme.
Exponentielle.
Puissances réelles.
Fonctions trigonométriques et réciproques.
Fonctions hyperboliques.
\end{tcolorbox}
\vfill
{\large Document original de TD avec corrigé complet}\par
{\large \today}\par
\end{titlepage}
\tableofcontents
\newpage
\section*{Avant-propos pédagogique}
\addcontentsline{toc}{section}{Avant-propos pédagogique}
Ce document constitue un TD complet de niveau classe préparatoire scientifique sur la dérivation des fonctions réelles d'une variable réelle.
Il a été construit pour prolonger un cours de type Analyse 4 portant sur la dérivabilité, le théorème de Rolle, le théorème des accroissements finis, l'inégalité des accroissements finis et les fonctions usuelles.
Les exercices sont originaux, progressifs, et rédigés pour entraîner à la fois le calcul, la preuve, la rédaction et les réflexes de concours.
\begin{attentionbox}
La ligne parfois recopiée dans des consignes, « applications linéaires, noyau, image, rang », relève de l'algèbre linéaire.
Elle n'appartient pas au chapitre d'analyse sur la dérivation.
Elle n'est donc pas utilisée dans ce TD afin de conserver une cohérence mathématique stricte.
\end{attentionbox}
\begin{methodebox}
Dans une copie de concours, une dérivée ne se calcule jamais « hors sol ».
On précise d'abord le domaine de définition et les intervalles de dérivabilité.
On identifie ensuite la structure de la fonction : somme, produit, quotient, composée ou réciproque.
Enfin, on exploite le signe de la dérivée avec le théorème des accroissements finis, et non par simple lecture graphique.
\end{methodebox}
\newpage

\section{Rappels méthodologiques indispensables}
\begin{methodebox}
\textbf{Définition locale.} $f$ est dérivable en $a$ si $\lim_{h\to0}\frac{f(a+h)-f(a)}h$ existe et est finie.
\end{methodebox}

\begin{methodebox}
\textbf{Approximation affine.} $f(a+h)=f(a)+h f'(a)+o(h)$ lorsque $h\to0$.
\end{methodebox}

\begin{methodebox}
\textbf{Tangente.} La tangente en $a$ est $y=f(a)+f'(a)(x-a)$.
\end{methodebox}

\begin{methodebox}
\textbf{Produit.} $(fg)'=f'g+fg'$ sur les intervalles de dérivabilité.
\end{methodebox}

\begin{methodebox}
\textbf{Quotient.} $(f/g)'=(f'g-fg')/g^2$ lorsque $g$ ne s'annule pas.
\end{methodebox}

\begin{methodebox}
\textbf{Composée.} $(f\circ g)'=(f'\circ g)g'$.
\end{methodebox}

\begin{methodebox}
\textbf{Réciproque.} $(f^{-1})'(y)=1/f'(f^{-1}(y))$ lorsque $f'$ ne s'annule pas au point correspondant.
\end{methodebox}

\begin{methodebox}
\textbf{Rolle.} Si $f(a)=f(b)$, avec continuité sur $[a,b]$ et dérivabilité sur $]a,b[$, alors une dérivée s'annule entre $a$ et $b$.
\end{methodebox}

\begin{methodebox}
\textbf{TAF.} Une pente moyenne est une pente instantanée.
\end{methodebox}

\begin{methodebox}
\textbf{IAF.} Une dérivée bornée donne une fonction lipschitzienne.
\end{methodebox}

\newpage
\section{Énoncés des exercices}
\subsection{Exercices d'application}

\begin{exobox}{1}{Nombre dérivé par définition -- \etoiles{$\star$}}

\begin{objectifbox}
Travailler un réflexe précis du chapitre : domaine, dérivabilité, calcul, théorème ou rédaction.
\end{objectifbox}

\begin{enumerate}

\item Soit $a\in\R$ et $f(x)=x^3$.

\item Calculer le taux d'accroissement $\dfrac{f(a+h)-f(a)}{h}$ pour $h
eq0$.

\item En déduire que $f$ est dérivable en $a$ et donner $f'(a)$.

\item Retrouver l'approximation affine de $f$ au voisinage de $a$.

\item Écrire l'équation de la tangente au graphe en $a=2$.

\end{enumerate}

\end{exobox}


\begin{exobox}{2}{Dérivées latérales d'une valeur absolue décalée -- \etoiles{$\star$}}

\begin{objectifbox}
Travailler un réflexe précis du chapitre : domaine, dérivabilité, calcul, théorème ou rédaction.
\end{objectifbox}

\begin{enumerate}

\item On pose $f(x)=|x-1|$.

\item Étudier la continuité de $f$ en $1$.

\item Calculer les dérivées à gauche et à droite en $1$.

\item Conclure sur la dérivabilité en $1$.

\item Interpréter géométriquement le résultat.

\end{enumerate}

\end{exobox}


\begin{exobox}{3}{Fonction avec raccord polynomial -- \etoiles{$\star\star$}}

\begin{objectifbox}
Travailler un réflexe précis du chapitre : domaine, dérivabilité, calcul, théorème ou rédaction.
\end{objectifbox}

\begin{enumerate}

\item On définit $f$ par $f(x)=x^2$ si $x\leq1$ et $f(x)=ax+b$ si $x>1$.

\item Trouver les conditions sur $a,b$ pour que $f$ soit continue en $1$.

\item Trouver les conditions sur $a,b$ pour que $f$ soit dérivable en $1$.

\item Donner alors l'expression complète de $f$.

\item Vérifier le résultat à l'aide des taux d'accroissement latéraux.

\end{enumerate}

\end{exobox}


\begin{exobox}{4}{Tangente et position locale -- \etoiles{$\star\star$}}

\begin{objectifbox}
Travailler un réflexe précis du chapitre : domaine, dérivabilité, calcul, théorème ou rédaction.
\end{objectifbox}

\begin{enumerate}

\item Soit $f(x)=\sqrt{1+x}$ définie sur $]-1,+\infty[$.

\item Calculer $f'(x)$ sur son domaine de dérivabilité.

\item Donner la tangente en $0$.

\item Étudier le signe de $f(x)-(1+x/2)$ pour $x>-1$ à l'aide d'une factorisation ou d'une inégalité simple.

\item En déduire la position du graphe par rapport à sa tangente en $0$.

\end{enumerate}

\end{exobox}


\begin{exobox}{5}{Calcul élémentaire de dérivée -- \etoiles{$\star$}}

\begin{objectifbox}
Travailler un réflexe précis du chapitre : domaine, dérivabilité, calcul, théorème ou rédaction.
\end{objectifbox}

\begin{enumerate}

\item Déterminer le domaine de dérivabilité de $f(x)=\dfrac{x^2+3x+1}{x-2}$.

\item Calculer $f'(x)$.

\item Simplifier le numérateur de $f'(x)$.

\item Résoudre $f'(x)=0$.

\item Donner le signe de $f'$ sur les intervalles du domaine.

\end{enumerate}

\end{exobox}


\begin{exobox}{6}{Composition logarithmique -- \etoiles{$\star\star$}}

\begin{objectifbox}
Travailler un réflexe précis du chapitre : domaine, dérivabilité, calcul, théorème ou rédaction.
\end{objectifbox}

\begin{enumerate}

\item Soit $f(x)=\ln(2+x^2)$.

\item Justifier que $f$ est définie et dérivable sur $\R$.

\item Calculer $f'(x)$.

\item Déterminer les variations de $f$.

\item Identifier un minimum global.

\end{enumerate}

\end{exobox}


\begin{exobox}{7}{Exponentielle composée -- \etoiles{$\star\star$}}

\begin{objectifbox}
Travailler un réflexe précis du chapitre : domaine, dérivabilité, calcul, théorème ou rédaction.
\end{objectifbox}

\begin{enumerate}

\item Soit $f(x)=\exp(-x^2+4x)$.

\item Calculer $f'(x)$.

\item Déterminer le signe de $f'(x)$.

\item Étudier les variations de $f$.

\item Donner le maximum global de $f$.

\end{enumerate}

\end{exobox}


\begin{exobox}{8}{Dérivée de puissance réelle -- \etoiles{$\star\star$}}

\begin{objectifbox}
Travailler un réflexe précis du chapitre : domaine, dérivabilité, calcul, théorème ou rédaction.
\end{objectifbox}

\begin{enumerate}

\item On fixe $\alpha\in\R$ et on pose $f(x)=x^\alpha\ln x$ sur $]0,+\infty[$.

\item Justifier la dérivabilité sur $]0,+\infty[$.

\item Calculer $f'(x)$.

\item Factoriser la dérivée.

\item Discuter rapidement le cas $\alpha=0$.

\end{enumerate}

\end{exobox}


\begin{exobox}{9}{Arctangente composée -- \etoiles{$\star\star$}}

\begin{objectifbox}
Travailler un réflexe précis du chapitre : domaine, dérivabilité, calcul, théorème ou rédaction.
\end{objectifbox}

\begin{enumerate}

\item Soit $f(x)=\arctan\left(\dfrac{2x}{1-x^2}\right)$ sur un intervalle ne contenant pas $\pm1$.

\item Calculer la dérivée de $u(x)=\dfrac{2x}{1-x^2}$.

\item Calculer $1+u(x)^2$.

\item En déduire $f'(x)$.

\item Interpréter le résultat sur un intervalle où l'expression est définie.

\end{enumerate}

\end{exobox}


\begin{exobox}{10}{Fonctions hyperboliques -- \etoiles{$\star\star$}}

\begin{objectifbox}
Travailler un réflexe précis du chapitre : domaine, dérivabilité, calcul, théorème ou rédaction.
\end{objectifbox}

\begin{enumerate}

\item On pose $f(x)=\ln(\ch x)$.

\item Justifier que $f$ est définie sur $\R$.

\item Calculer $f'(x)$.

\item Montrer que $f$ est paire.

\item Déterminer ses variations sur $[0,+\infty[$.

\end{enumerate}

\end{exobox}


\begin{exobox}{11}{Dérivabilité d'une racine -- \etoiles{$\star\star$}}

\begin{objectifbox}
Travailler un réflexe précis du chapitre : domaine, dérivabilité, calcul, théorème ou rédaction.
\end{objectifbox}

\begin{enumerate}

\item Soit $f(x)=\sqrt{|x|}$.

\item Étudier la continuité en $0$.

\item Calculer le taux $\dfrac{f(h)-f(0)}{h}$ pour $h>0$ et pour $h<0$.

\item Décider si $f$ est dérivable en $0$.

\item Comparer avec la fonction $x\mapsto |x|$.

\end{enumerate}

\end{exobox}


\begin{exobox}{12}{Dérivée d'une réciproque simple -- \etoiles{$\star\star$}}

\begin{objectifbox}
Travailler un réflexe précis du chapitre : domaine, dérivabilité, calcul, théorème ou rédaction.
\end{objectifbox}

\begin{enumerate}

\item Soit $f(x)=x+x^3$.

\item Montrer que $f$ réalise une bijection de $\R$ sur $\R$.

\item Justifier que $f^{-1}$ est dérivable sur $\R$.

\item Donner une formule pour $(f^{-1})'(y)$.

\item Calculer $(f^{-1})'(0)$.

\end{enumerate}

\end{exobox}


\subsection{Exercices de méthode}

\begin{exobox}{13}{Étude complète d'un quotient -- \etoiles{$\star\star\star$}}

\begin{objectifbox}
Travailler un réflexe précis du chapitre : domaine, dérivabilité, calcul, théorème ou rédaction.
\end{objectifbox}

\begin{enumerate}

\item On considère la fonction $f(x)=\frac{x^2+1}{x^2-1}$.

\item Déterminer soigneusement le domaine de définition et les intervalles de dérivabilité.

\item Calculer la dérivée et vérifier que l'expression attendue est $f'(x)=\frac{-4x}{(x^2-1)^2}$.

\item Étudier le signe de la dérivée lorsque cela est possible sans méthode numérique.

\item En déduire un tableau de variations ou une propriété qualitative précise.

\end{enumerate}

\end{exobox}


\begin{exobox}{14}{Logarithme et variations -- \etoiles{$\star\star\star$}}

\begin{objectifbox}
Travailler un réflexe précis du chapitre : domaine, dérivabilité, calcul, théorème ou rédaction.
\end{objectifbox}

\begin{enumerate}

\item On considère la fonction $f(x)=x\ln x-x$.

\item Déterminer soigneusement le domaine de définition et les intervalles de dérivabilité.

\item Calculer la dérivée et vérifier que l'expression attendue est $f'(x)=\ln x$.

\item Étudier le signe de la dérivée lorsque cela est possible sans méthode numérique.

\item En déduire un tableau de variations ou une propriété qualitative précise.

\end{enumerate}

\end{exobox}


\begin{exobox}{15}{Inégalité classique par étude de fonction -- \etoiles{$\star\star\star$}}

\begin{objectifbox}
Travailler un réflexe précis du chapitre : domaine, dérivabilité, calcul, théorème ou rédaction.
\end{objectifbox}

\begin{enumerate}

\item On considère la fonction $f(x)=\e^x-1-x$.

\item Déterminer soigneusement le domaine de définition et les intervalles de dérivabilité.

\item Calculer la dérivée et vérifier que l'expression attendue est $f'(x)=\e^x-1$.

\item Étudier le signe de la dérivée lorsque cela est possible sans méthode numérique.

\item En déduire un tableau de variations ou une propriété qualitative précise.

\end{enumerate}

\end{exobox}


\begin{exobox}{16}{Puissance et logarithme -- \etoiles{$\star\star\star$}}

\begin{objectifbox}
Travailler un réflexe précis du chapitre : domaine, dérivabilité, calcul, théorème ou rédaction.
\end{objectifbox}

\begin{enumerate}

\item On considère la fonction $f(x)=\frac{\ln x}{x}$.

\item Déterminer soigneusement le domaine de définition et les intervalles de dérivabilité.

\item Calculer la dérivée et vérifier que l'expression attendue est $f'(x)=\frac{1-\ln x}{x^2}$.

\item Étudier le signe de la dérivée lorsque cela est possible sans méthode numérique.

\item En déduire un tableau de variations ou une propriété qualitative précise.

\end{enumerate}

\end{exobox}


\begin{exobox}{17}{Tangente horizontale et extrema -- \etoiles{$\star\star\star$}}

\begin{objectifbox}
Travailler un réflexe précis du chapitre : domaine, dérivabilité, calcul, théorème ou rédaction.
\end{objectifbox}

\begin{enumerate}

\item On considère la fonction $f(x)=x^3-3x+1$.

\item Déterminer soigneusement le domaine de définition et les intervalles de dérivabilité.

\item Calculer la dérivée et vérifier que l'expression attendue est $f'(x)=3x^2-3$.

\item Étudier le signe de la dérivée lorsque cela est possible sans méthode numérique.

\item En déduire un tableau de variations ou une propriété qualitative précise.

\end{enumerate}

\end{exobox}


\begin{exobox}{18}{Trigonométrie composée -- \etoiles{$\star\star\star$}}

\begin{objectifbox}
Travailler un réflexe précis du chapitre : domaine, dérivabilité, calcul, théorème ou rédaction.
\end{objectifbox}

\begin{enumerate}

\item On considère la fonction $f(x)=\sin x+\cos(2x)$.

\item Déterminer soigneusement le domaine de définition et les intervalles de dérivabilité.

\item Calculer la dérivée et vérifier que l'expression attendue est $f'(x)=\cos x-2\sin(2x)$.

\item Étudier le signe de la dérivée lorsque cela est possible sans méthode numérique.

\item En déduire un tableau de variations ou une propriété qualitative précise.

\end{enumerate}

\end{exobox}


\begin{exobox}{19}{Arcsinus et domaine -- \etoiles{$\star\star\star$}}

\begin{objectifbox}
Travailler un réflexe précis du chapitre : domaine, dérivabilité, calcul, théorème ou rédaction.
\end{objectifbox}

\begin{enumerate}

\item On considère la fonction $f(x)=\arcsin(2x-1)$.

\item Déterminer soigneusement le domaine de définition et les intervalles de dérivabilité.

\item Calculer la dérivée et vérifier que l'expression attendue est $f'(x)=\frac{1}{\sqrt{x(1-x)}}\text{ sur }]0,1[$.

\item Étudier le signe de la dérivée lorsque cela est possible sans méthode numérique.

\item En déduire un tableau de variations ou une propriété qualitative précise.

\end{enumerate}

\end{exobox}


\begin{exobox}{20}{Arccosinus et monotonie -- \etoiles{$\star\star\star$}}

\begin{objectifbox}
Travailler un réflexe précis du chapitre : domaine, dérivabilité, calcul, théorème ou rédaction.
\end{objectifbox}

\begin{enumerate}

\item On considère la fonction $f(x)=\arccos\left(\frac{x}{\sqrt{1+x^2}}\right)$.

\item Déterminer soigneusement le domaine de définition et les intervalles de dérivabilité.

\item Calculer la dérivée et vérifier que l'expression attendue est $f'(x)=-\frac{1}{1+x^2}$.

\item Étudier le signe de la dérivée lorsque cela est possible sans méthode numérique.

\item En déduire un tableau de variations ou une propriété qualitative précise.

\end{enumerate}

\end{exobox}


\begin{exobox}{21}{Exponentielle rationnelle -- \etoiles{$\star\star\star$}}

\begin{objectifbox}
Travailler un réflexe précis du chapitre : domaine, dérivabilité, calcul, théorème ou rédaction.
\end{objectifbox}

\begin{enumerate}

\item On considère la fonction $f(x)=\e^{x}/(1+\e^{x})$.

\item Déterminer soigneusement le domaine de définition et les intervalles de dérivabilité.

\item Calculer la dérivée et vérifier que l'expression attendue est $f'(x)=\frac{\e^x}{(1+\e^x)^2}$.

\item Étudier le signe de la dérivée lorsque cela est possible sans méthode numérique.

\item En déduire un tableau de variations ou une propriété qualitative précise.

\end{enumerate}

\end{exobox}


\begin{exobox}{22}{Fonction hyperbolique quotient -- \etoiles{$\star\star\star$}}

\begin{objectifbox}
Travailler un réflexe précis du chapitre : domaine, dérivabilité, calcul, théorème ou rédaction.
\end{objectifbox}

\begin{enumerate}

\item On considère la fonction $f(x)=\thh x$.

\item Déterminer soigneusement le domaine de définition et les intervalles de dérivabilité.

\item Calculer la dérivée et vérifier que l'expression attendue est $f'(x)=\frac{1}{\ch^2x}$.

\item Étudier le signe de la dérivée lorsque cela est possible sans méthode numérique.

\item En déduire un tableau de variations ou une propriété qualitative précise.

\end{enumerate}

\end{exobox}


\begin{exobox}{23}{Dérivée logarithmique -- \etoiles{$\star\star\star$}}

\begin{objectifbox}
Travailler un réflexe précis du chapitre : domaine, dérivabilité, calcul, théorème ou rédaction.
\end{objectifbox}

\begin{enumerate}

\item On considère la fonction $f(x)=(x^2+1)^x$.

\item Déterminer soigneusement le domaine de définition et les intervalles de dérivabilité.

\item Calculer la dérivée et vérifier que l'expression attendue est $f'(x)=(x^2+1)^x\left(\ln(x^2+1)+\frac{2x^2}{x^2+1}\right)$.

\item Étudier le signe de la dérivée lorsque cela est possible sans méthode numérique.

\item En déduire un tableau de variations ou une propriété qualitative précise.

\end{enumerate}

\end{exobox}


\begin{exobox}{24}{Optimisation élémentaire -- \etoiles{$\star\star\star$}}

\begin{objectifbox}
Travailler un réflexe précis du chapitre : domaine, dérivabilité, calcul, théorème ou rédaction.
\end{objectifbox}

\begin{enumerate}

\item On considère la fonction $f(x)=x^a(1-x)^b$.

\item Déterminer soigneusement le domaine de définition et les intervalles de dérivabilité.

\item Calculer la dérivée et vérifier que l'expression attendue est $f'(x)=x^a(1-x)^b\left(\frac{a}{x}-\frac{b}{1-x}\right)$.

\item Étudier le signe de la dérivée lorsque cela est possible sans méthode numérique.

\item En déduire un tableau de variations ou une propriété qualitative précise.

\end{enumerate}

\end{exobox}


\subsection{Exercices de démonstration}

\begin{exobox}{25}{Dérivabilité implique continuité -- \etoiles{$\star\star\star$}}

\begin{objectifbox}
Travailler un réflexe précis du chapitre : domaine, dérivabilité, calcul, théorème ou rédaction.
\end{objectifbox}

\begin{enumerate}

\item Soit $f:I\to\R$ dérivable en $a\in I$.

\item Partir de la définition du nombre dérivé.

\item Construire une fonction $\varepsilon(h)$ tendant vers $0$.

\item Démontrer que $f(a+h)\to f(a)$.

\item Expliquer pourquoi la réciproque est fausse.

\end{enumerate}

\end{exobox}


\begin{exobox}{26}{Formule du produit par les taux -- \etoiles{$\star\star\star$}}

\begin{objectifbox}
Travailler un réflexe précis du chapitre : domaine, dérivabilité, calcul, théorème ou rédaction.
\end{objectifbox}

\begin{enumerate}

\item Soient $f$ et $g$ dérivables en $a$.

\item Écrire $(fg)(a+h)-(fg)(a)$ en ajoutant et retranchant $f(a)g(a+h)$ ou $f(a+h)g(a)$.

\item Diviser par $h$.

\item Utiliser la continuité de $f$ ou de $g$.

\item Conclure la formule $(fg)'(a)=f'(a)g(a)+f(a)g'(a)$.

\end{enumerate}

\end{exobox}


\begin{exobox}{27}{Formule de la composée -- \etoiles{$\star\star\star\star$}}

\begin{objectifbox}
Travailler un réflexe précis du chapitre : domaine, dérivabilité, calcul, théorème ou rédaction.
\end{objectifbox}

\begin{enumerate}

\item Soient $g$ dérivable en $a$ et $f$ dérivable en $g(a)$.

\item Écrire le développement d'ordre $1$ de $f$ en $g(a)$.

\item Poser $k=g(a+h)-g(a)$.

\item Utiliser la dérivabilité de $g$ en $a$.

\item Conclure $(f\circ g)'(a)=f'(g(a))g'(a)$.

\end{enumerate}

\end{exobox}


\begin{exobox}{28}{Dérivée de la réciproque -- \etoiles{$\star\star\star\star$}}

\begin{objectifbox}
Travailler un réflexe précis du chapitre : domaine, dérivabilité, calcul, théorème ou rédaction.
\end{objectifbox}

\begin{enumerate}

\item Soit $f:I\to J$ une bijection continue strictement monotone, dérivable en $a$, avec $f'(a)
e0$.

\item Poser $b=f(a)$.

\item Écrire le taux d'accroissement de $f^{-1}$ en $b$.

\item Faire le changement de variable $y=f(x)$.

\item Conclure la formule de dérivation de la réciproque.

\end{enumerate}

\end{exobox}


\begin{exobox}{29}{Condition nécessaire d'extremum -- \etoiles{$\star\star\star$}}

\begin{objectifbox}
Travailler un réflexe précis du chapitre : domaine, dérivabilité, calcul, théorème ou rédaction.
\end{objectifbox}

\begin{enumerate}

\item Soit $f$ dérivable en un point intérieur $a$ d'un intervalle $I$.

\item On suppose que $f$ admet un maximum local en $a$.

\item Étudier les quotients pour $h>0$ et pour $h<0$.

\item Passer à la limite.

\item Conclure et traiter brièvement le cas du minimum.

\end{enumerate}

\end{exobox}


\begin{exobox}{30}{Théorème de Rolle -- \etoiles{$\star\star\star\star$}}

\begin{objectifbox}
Travailler un réflexe précis du chapitre : domaine, dérivabilité, calcul, théorème ou rédaction.
\end{objectifbox}

\begin{enumerate}

\item Énoncer précisément le théorème de Rolle.

\item Démontrer le résultat à l'aide du théorème des bornes atteintes.

\item Distinguer le cas où $f$ est constante.

\item Dans le cas non constant, expliquer pourquoi un extremum est atteint à l'intérieur.

\item Conclure par la condition nécessaire d'extremum local.

\end{enumerate}

\end{exobox}


\begin{exobox}{31}{Théorème des accroissements finis -- \etoiles{$\star\star\star\star$}}

\begin{objectifbox}
Travailler un réflexe précis du chapitre : domaine, dérivabilité, calcul, théorème ou rédaction.
\end{objectifbox}

\begin{enumerate}

\item Soit $f$ continue sur $[a,b]$ et dérivable sur $]a,b[$.

\item Construire une fonction auxiliaire $g$ à laquelle appliquer Rolle.

\item Vérifier $g(a)=g(b)$.

\item Calculer $g'$.

\item Conclure le théorème des accroissements finis.

\end{enumerate}

\end{exobox}


\begin{exobox}{32}{Inégalité des accroissements finis -- \etoiles{$\star\star\star\star$}}

\begin{objectifbox}
Travailler un réflexe précis du chapitre : domaine, dérivabilité, calcul, théorème ou rédaction.
\end{objectifbox}

\begin{enumerate}

\item Soit $f$ dérivable sur un intervalle $I$ et $|f'|\le K$ sur $I$.

\item Prendre $x<y$ dans $I$.

\item Appliquer le théorème des accroissements finis sur $[x,y]$.

\item Majorer la valeur absolue obtenue.

\item Conclure pour tous $x,y$.

\end{enumerate}

\end{exobox}


\begin{exobox}{33}{Rolle et racines successives -- \etoiles{$\star\star\star\star$}}

\begin{objectifbox}
Travailler un réflexe précis du chapitre : domaine, dérivabilité, calcul, théorème ou rédaction.
\end{objectifbox}

\begin{enumerate}

\item Soit $P$ un polynôme réel ayant $n$ racines réelles distinctes.

\item Montrer que $P'$ possède au moins $n-1$ racines réelles distinctes.

\item Préciser où se trouvent ces racines.

\item Appliquer le résultat à un polynôme de degré $n$ scindé à racines simples.

\item Donner une conséquence qualitative.

\end{enumerate}

\end{exobox}


\begin{exobox}{34}{Convexité sans nom par dérivée croissante -- \etoiles{$\star\star\star\star$}}

\begin{objectifbox}
Travailler un réflexe précis du chapitre : domaine, dérivabilité, calcul, théorème ou rédaction.
\end{objectifbox}

\begin{enumerate}

\item Soit $f$ dérivable sur un intervalle $I$ et supposons $f'$ croissante.

\item Pour $a<x<b$, appliquer le TAF à $[a,x]$ et $[x,b]$.

\item Comparer les pentes moyennes.

\item En déduire une inégalité reliant $f(x)$ à $f(a)$ et $f(b)$.

\item Ne pas employer le vocabulaire de convexité si le cours ne l'a pas encore introduit.

\end{enumerate}

\end{exobox}


\subsection{Exercices de réflexion et préparation concours}

\begin{exobox}{35}{Inégalité $\ln x\le x-1$ et raffinements -- \etoiles{$\star\star\star\star$}}

\begin{objectifbox}
Travailler un réflexe précis du chapitre : domaine, dérivabilité, calcul, théorème ou rédaction.
\end{objectifbox}

\begin{enumerate}

\item On considère $F(x)=x-1-\ln x$ sur son domaine naturel.

\item Déterminer le domaine et les intervalles de dérivabilité.

\item Calculer la dérivée première.

\item Étudier le signe de la dérivée en justifiant chaque changement de signe.

\item En déduire les variations et les extrema éventuels.

\item Utiliser l'étude pour prouver une inégalité ou une unicité d'équation.

\end{enumerate}

\end{exobox}


\begin{exobox}{36}{Encadrement de l'exponentielle -- \etoiles{$\star\star\star\star$}}

\begin{objectifbox}
Travailler un réflexe précis du chapitre : domaine, dérivabilité, calcul, théorème ou rédaction.
\end{objectifbox}

\begin{enumerate}

\item On considère $F(x)=\e^x-1-x$ sur son domaine naturel.

\item Déterminer le domaine et les intervalles de dérivabilité.

\item Calculer la dérivée première.

\item Étudier le signe de la dérivée en justifiant chaque changement de signe.

\item En déduire les variations et les extrema éventuels.

\item Utiliser l'étude pour prouver une inégalité ou une unicité d'équation.

\end{enumerate}

\end{exobox}


\begin{exobox}{37}{Unicité pour une équation transcendante -- \etoiles{$\star\star\star\star$}}

\begin{objectifbox}
Travailler un réflexe précis du chapitre : domaine, dérivabilité, calcul, théorème ou rédaction.
\end{objectifbox}

\begin{enumerate}

\item On considère $F(x)=x+\ln x$ sur son domaine naturel.

\item Déterminer le domaine et les intervalles de dérivabilité.

\item Calculer la dérivée première.

\item Étudier le signe de la dérivée en justifiant chaque changement de signe.

\item En déduire les variations et les extrema éventuels.

\item Utiliser l'étude pour prouver une inégalité ou une unicité d'équation.

\end{enumerate}

\end{exobox}


\begin{exobox}{38}{Point fixe de $\cos$ -- \etoiles{$\star\star\star\star$}}

\begin{objectifbox}
Travailler un réflexe précis du chapitre : domaine, dérivabilité, calcul, théorème ou rédaction.
\end{objectifbox}

\begin{enumerate}

\item On considère $F(x)=\cos x-x$ sur son domaine naturel.

\item Déterminer le domaine et les intervalles de dérivabilité.

\item Calculer la dérivée première.

\item Étudier le signe de la dérivée en justifiant chaque changement de signe.

\item En déduire les variations et les extrema éventuels.

\item Utiliser l'étude pour prouver une inégalité ou une unicité d'équation.

\end{enumerate}

\end{exobox}


\begin{exobox}{39}{Étude de $x-\arctan x$ -- \etoiles{$\star\star\star\star$}}

\begin{objectifbox}
Travailler un réflexe précis du chapitre : domaine, dérivabilité, calcul, théorème ou rédaction.
\end{objectifbox}

\begin{enumerate}

\item On considère $F(x)=x-\arctan x$ sur son domaine naturel.

\item Déterminer le domaine et les intervalles de dérivabilité.

\item Calculer la dérivée première.

\item Étudier le signe de la dérivée en justifiant chaque changement de signe.

\item En déduire les variations et les extrema éventuels.

\item Utiliser l'étude pour prouver une inégalité ou une unicité d'équation.

\end{enumerate}

\end{exobox}


\begin{exobox}{40}{Étude de $\sin x/x$ -- \etoiles{$\star\star\star\star$}}

\begin{objectifbox}
Travailler un réflexe précis du chapitre : domaine, dérivabilité, calcul, théorème ou rédaction.
\end{objectifbox}

\begin{enumerate}

\item On considère $F(x)=\frac{\sin x}{x}$ sur son domaine naturel.

\item Déterminer le domaine et les intervalles de dérivabilité.

\item Calculer la dérivée première.

\item Étudier le signe de la dérivée en justifiant chaque changement de signe.

\item En déduire les variations et les extrema éventuels.

\item Utiliser l'étude pour prouver une inégalité ou une unicité d'équation.

\end{enumerate}

\end{exobox}


\begin{exobox}{41}{Maximum d'une fonction avec paramètre -- \etoiles{$\star\star\star\star$}}

\begin{objectifbox}
Travailler un réflexe précis du chapitre : domaine, dérivabilité, calcul, théorème ou rédaction.
\end{objectifbox}

\begin{enumerate}

\item On considère $F_a(x)=x^a\e^{-x}$ sur son domaine naturel.

\item Déterminer le domaine et les intervalles de dérivabilité.

\item Calculer la dérivée première.

\item Étudier le signe de la dérivée en justifiant chaque changement de signe.

\item En déduire les variations et les extrema éventuels.

\item Utiliser l'étude pour prouver une inégalité ou une unicité d'équation.

\end{enumerate}

\end{exobox}


\begin{exobox}{42}{Réciproque de $x+\e^x$ -- \etoiles{$\star\star\star\star$}}

\begin{objectifbox}
Travailler un réflexe précis du chapitre : domaine, dérivabilité, calcul, théorème ou rédaction.
\end{objectifbox}

\begin{enumerate}

\item On considère $F(x)=x+\e^x$ sur son domaine naturel.

\item Déterminer le domaine et les intervalles de dérivabilité.

\item Calculer la dérivée première.

\item Étudier le signe de la dérivée en justifiant chaque changement de signe.

\item En déduire les variations et les extrema éventuels.

\item Utiliser l'étude pour prouver une inégalité ou une unicité d'équation.

\end{enumerate}

\end{exobox}


\begin{exobox}{43}{Comparaison logarithme-puissance -- \etoiles{$\star\star\star\star$}}

\begin{objectifbox}
Travailler un réflexe précis du chapitre : domaine, dérivabilité, calcul, théorème ou rédaction.
\end{objectifbox}

\begin{enumerate}

\item On considère $F_\alpha(x)=\frac{\ln x}{x^\alpha}$ sur son domaine naturel.

\item Déterminer le domaine et les intervalles de dérivabilité.

\item Calculer la dérivée première.

\item Étudier le signe de la dérivée en justifiant chaque changement de signe.

\item En déduire les variations et les extrema éventuels.

\item Utiliser l'étude pour prouver une inégalité ou une unicité d'équation.

\end{enumerate}

\end{exobox}


\begin{exobox}{44}{Hyperboliques et inégalités -- \etoiles{$\star\star\star\star$}}

\begin{objectifbox}
Travailler un réflexe précis du chapitre : domaine, dérivabilité, calcul, théorème ou rédaction.
\end{objectifbox}

\begin{enumerate}

\item On considère $F(x)=\ch x-1-\frac{x^2}{2}$ sur son domaine naturel.

\item Déterminer le domaine et les intervalles de dérivabilité.

\item Calculer la dérivée première.

\item Étudier le signe de la dérivée en justifiant chaque changement de signe.

\item En déduire les variations et les extrema éventuels.

\item Utiliser l'étude pour prouver une inégalité ou une unicité d'équation.

\end{enumerate}

\end{exobox}


\subsection{Exercices type oraux Mines, Centrale, CCP}

\begin{exobox}{45}{Problème CCP -- Rolle itéré et polynômes -- \etoiles{$\star\star\star\star\star$}}

\begin{objectifbox}
Travailler un réflexe précis du chapitre : domaine, dérivabilité, calcul, théorème ou rédaction.
\end{objectifbox}

\begin{enumerate}

\item Ce problème est conçu comme un enchaînement de questions de concours.

\item On demande une rédaction complète et non une simple suite de calculs.

\item Question A. Établir le domaine, la continuité et la dérivabilité sur les intervalles pertinents.

\item Question B. Calculer la dérivée ou les dérivées utiles, puis factoriser le résultat.

\item Question C. Appliquer Rolle, le TAF ou l'IAF pour obtenir une propriété globale.

\item Question D. Résoudre une équation ou démontrer une inégalité précise.

\item Question E. Conclure par une synthèse claire indiquant existence, unicité, variation ou borne optimale.

\end{enumerate}

\end{exobox}


\begin{exobox}{46}{Problème Mines -- Inégalités par TAF -- \etoiles{$\star\star\star\star\star$}}

\begin{objectifbox}
Travailler un réflexe précis du chapitre : domaine, dérivabilité, calcul, théorème ou rédaction.
\end{objectifbox}

\begin{enumerate}

\item Ce problème est conçu comme un enchaînement de questions de concours.

\item On demande une rédaction complète et non une simple suite de calculs.

\item Question A. Établir le domaine, la continuité et la dérivabilité sur les intervalles pertinents.

\item Question B. Calculer la dérivée ou les dérivées utiles, puis factoriser le résultat.

\item Question C. Appliquer Rolle, le TAF ou l'IAF pour obtenir une propriété globale.

\item Question D. Résoudre une équation ou démontrer une inégalité précise.

\item Question E. Conclure par une synthèse claire indiquant existence, unicité, variation ou borne optimale.

\end{enumerate}

\end{exobox}


\begin{exobox}{47}{Problème Centrale -- Réciproque et paramètre -- \etoiles{$\star\star\star\star\star$}}

\begin{objectifbox}
Travailler un réflexe précis du chapitre : domaine, dérivabilité, calcul, théorème ou rédaction.
\end{objectifbox}

\begin{enumerate}

\item Ce problème est conçu comme un enchaînement de questions de concours.

\item On demande une rédaction complète et non une simple suite de calculs.

\item Question A. Établir le domaine, la continuité et la dérivabilité sur les intervalles pertinents.

\item Question B. Calculer la dérivée ou les dérivées utiles, puis factoriser le résultat.

\item Question C. Appliquer Rolle, le TAF ou l'IAF pour obtenir une propriété globale.

\item Question D. Résoudre une équation ou démontrer une inégalité précise.

\item Question E. Conclure par une synthèse claire indiquant existence, unicité, variation ou borne optimale.

\end{enumerate}

\end{exobox}


\begin{exobox}{48}{Problème oral -- Tangentes communes -- \etoiles{$\star\star\star\star\star$}}

\begin{objectifbox}
Travailler un réflexe précis du chapitre : domaine, dérivabilité, calcul, théorème ou rédaction.
\end{objectifbox}

\begin{enumerate}

\item Ce problème est conçu comme un enchaînement de questions de concours.

\item On demande une rédaction complète et non une simple suite de calculs.

\item Question A. Établir le domaine, la continuité et la dérivabilité sur les intervalles pertinents.

\item Question B. Calculer la dérivée ou les dérivées utiles, puis factoriser le résultat.

\item Question C. Appliquer Rolle, le TAF ou l'IAF pour obtenir une propriété globale.

\item Question D. Résoudre une équation ou démontrer une inégalité précise.

\item Question E. Conclure par une synthèse claire indiquant existence, unicité, variation ou borne optimale.

\end{enumerate}

\end{exobox}


\begin{exobox}{49}{Problème oral -- Arctangente et logarithme -- \etoiles{$\star\star\star\star\star$}}

\begin{objectifbox}
Travailler un réflexe précis du chapitre : domaine, dérivabilité, calcul, théorème ou rédaction.
\end{objectifbox}

\begin{enumerate}

\item Ce problème est conçu comme un enchaînement de questions de concours.

\item On demande une rédaction complète et non une simple suite de calculs.

\item Question A. Établir le domaine, la continuité et la dérivabilité sur les intervalles pertinents.

\item Question B. Calculer la dérivée ou les dérivées utiles, puis factoriser le résultat.

\item Question C. Appliquer Rolle, le TAF ou l'IAF pour obtenir une propriété globale.

\item Question D. Résoudre une équation ou démontrer une inégalité précise.

\item Question E. Conclure par une synthèse claire indiquant existence, unicité, variation ou borne optimale.

\end{enumerate}

\end{exobox}


\begin{exobox}{50}{Problème bilan -- Fonction définie par raccords -- \etoiles{$\star\star\star\star\star$}}

\begin{objectifbox}
Travailler un réflexe précis du chapitre : domaine, dérivabilité, calcul, théorème ou rédaction.
\end{objectifbox}

\begin{enumerate}

\item Ce problème est conçu comme un enchaînement de questions de concours.

\item On demande une rédaction complète et non une simple suite de calculs.

\item Question A. Établir le domaine, la continuité et la dérivabilité sur les intervalles pertinents.

\item Question B. Calculer la dérivée ou les dérivées utiles, puis factoriser le résultat.

\item Question C. Appliquer Rolle, le TAF ou l'IAF pour obtenir une propriété globale.

\item Question D. Résoudre une équation ou démontrer une inégalité précise.

\item Question E. Conclure par une synthèse claire indiquant existence, unicité, variation ou borne optimale.

\end{enumerate}

\end{exobox}


\newpage
\section{Corrigé complet et détaillé}
\begin{corrbox}{1}{Nombre dérivé par définition}

\begin{methodebox}
On rédige en vérifiant d'abord les hypothèses : domaine, continuité, dérivabilité, puis calcul ou théorème utilisé.
\end{methodebox}

Pour $h
eq0$, on calcule $(a+h)^3-a^3=3a^2h+3ah^2+h^3$.


Ainsi le taux d'accroissement vaut $3a^2+3ah+h^2$.


Lorsque $h\to0$, ce quotient tend vers $3a^2$.


La fonction est donc dérivable en $a$ et $f'(a)=3a^2$.


L'approximation affine est $f(a+h)=a^3+3a^2h+o(h)$.


En $a=2$, on a $f(2)=8$ et $f'(2)=12$, donc la tangente est $y=8+12(x-2)$, c'est-à-dire $y=12x-16$.


\begin{attentionbox}
Vérification finale : la conclusion doit porter sur l'intervalle exact de travail et ne doit pas franchir un point exclu du domaine.
\end{attentionbox}

\end{corrbox}


\begin{corrbox}{2}{Dérivées latérales d'une valeur absolue décalée}

\begin{methodebox}
On rédige en vérifiant d'abord les hypothèses : domaine, continuité, dérivabilité, puis calcul ou théorème utilisé.
\end{methodebox}

La fonction valeur absolue est continue sur $\R$, donc $f$ est continue en $1$.


Pour $h>0$, $f(1+h)=h$, donc $\frac{f(1+h)-f(1)}{h}=1$.


Pour $h<0$, $f(1+h)=|-h|=-h$, donc le quotient vaut $-1$.


La dérivée à droite vaut $1$ et la dérivée à gauche vaut $-1$.


Ces deux limites étant différentes, $f$ n'est pas dérivable en $1$.


Le graphe présente un angle au point d'abscisse $1$.


\begin{attentionbox}
Vérification finale : la conclusion doit porter sur l'intervalle exact de travail et ne doit pas franchir un point exclu du domaine.
\end{attentionbox}

\end{corrbox}


\begin{corrbox}{3}{Fonction avec raccord polynomial}

\begin{methodebox}
On rédige en vérifiant d'abord les hypothèses : domaine, continuité, dérivabilité, puis calcul ou théorème utilisé.
\end{methodebox}

La continuité en $1$ impose $\lim_{x\to1^-}x^2=f(1)=1$ et $\lim_{x\to1^+}(ax+b)=a+b$.


Il faut donc $a+b=1$.


La dérivée à gauche de $x^2$ en $1$ vaut $2$.


La dérivée à droite de $ax+b$ vaut $a$.


La dérivabilité impose donc $a=2$.


Avec $a+b=1$, on obtient $b=-1$.


La fonction dérivable est donc $f(x)=x^2$ si $x\le1$ et $f(x)=2x-1$ si $x>1$.


Les deux quotients latéraux tendent bien vers $2$.


\begin{attentionbox}
Vérification finale : la conclusion doit porter sur l'intervalle exact de travail et ne doit pas franchir un point exclu du domaine.
\end{attentionbox}

\end{corrbox}


\begin{corrbox}{4}{Tangente et position locale}

\begin{methodebox}
On rédige en vérifiant d'abord les hypothèses : domaine, continuité, dérivabilité, puis calcul ou théorème utilisé.
\end{methodebox}

On écrit $f(x)=(1+x)^{1/2}$.


Par dérivation d'une puissance composée, $f'(x)=\frac1{2\sqrt{1+x}}$.


En $0$, $f(0)=1$ et $f'(0)=\frac12$.


La tangente est donc $y=1+\frac{x}{2}$.


On compare $\sqrt{1+x}$ et $1+x/2$.


Pour $x>-1$, l'inégalité $\sqrt{1+x}\le 1+x/2$ équivaut, lorsque le membre de droite est positif, à $1+x\le 1+x+x^2/4$.


Cette dernière inégalité est vraie, avec égalité seulement pour $x=0$.


Le graphe est donc en dessous de sa tangente en $0$.


\begin{attentionbox}
Vérification finale : la conclusion doit porter sur l'intervalle exact de travail et ne doit pas franchir un point exclu du domaine.
\end{attentionbox}

\end{corrbox}


\begin{corrbox}{5}{Calcul élémentaire de dérivée}

\begin{methodebox}
On rédige en vérifiant d'abord les hypothèses : domaine, continuité, dérivabilité, puis calcul ou théorème utilisé.
\end{methodebox}

La fonction est définie et dérivable sur $\R\setminus\{2\}$.


On pose $u=x^2+3x+1$ et $v=x-2$.


Alors $u'=2x+3$ et $v'=1$.


La formule du quotient donne $f'=\frac{(2x+3)(x-2)-(x^2+3x+1)}{(x-2)^2}$.


Le numérateur vaut $2x^2-x-6-x^2-3x-1=x^2-4x-7$.


Donc $f'(x)=\frac{x^2-4x-7}{(x-2)^2}$.


Les zéros sont $x=2\pm\sqrt{11}$.


Le signe de $f'$ est celui du trinôme $x^2-4x-7$, car le dénominateur est strictement positif sur le domaine.


\begin{attentionbox}
Vérification finale : la conclusion doit porter sur l'intervalle exact de travail et ne doit pas franchir un point exclu du domaine.
\end{attentionbox}

\end{corrbox}


\begin{corrbox}{6}{Composition logarithmique}

\begin{methodebox}
On rédige en vérifiant d'abord les hypothèses : domaine, continuité, dérivabilité, puis calcul ou théorème utilisé.
\end{methodebox}

Pour tout $x\in\R$, $2+x^2>0$, donc le logarithme est bien défini.


La fonction est composée de fonctions dérivables, donc dérivable sur $\R$.


On pose $u(x)=2+x^2$, d'où $u'(x)=2x$.


Ainsi $f'(x)=\frac{2x}{2+x^2}$.


Le dénominateur est positif, donc le signe de $f'$ est celui de $x$.


La fonction est décroissante sur $]-\infty,0]$ puis croissante sur $[0,+\infty[$.


Elle admet un minimum global en $0$, de valeur $\ln2$.


\begin{attentionbox}
Vérification finale : la conclusion doit porter sur l'intervalle exact de travail et ne doit pas franchir un point exclu du domaine.
\end{attentionbox}

\end{corrbox}


\begin{corrbox}{7}{Exponentielle composée}

\begin{methodebox}
On rédige en vérifiant d'abord les hypothèses : domaine, continuité, dérivabilité, puis calcul ou théorème utilisé.
\end{methodebox}

On pose $u(x)=-x^2+4x$.


Alors $u'(x)=-2x+4=2(2-x)$.


La dérivée de $\e^{u(x)}$ est $u'(x)\e^{u(x)}$.


Ainsi $f'(x)=2(2-x)\e^{-x^2+4x}$.


L'exponentielle est strictement positive, donc le signe est celui de $2-x$.


La fonction est croissante sur $]-\infty,2]$ puis décroissante sur $[2,+\infty[$.


Le maximum global vaut $f(2)=\e^4$.


\begin{attentionbox}
Vérification finale : la conclusion doit porter sur l'intervalle exact de travail et ne doit pas franchir un point exclu du domaine.
\end{attentionbox}

\end{corrbox}


\begin{corrbox}{8}{Dérivée de puissance réelle}

\begin{methodebox}
On rédige en vérifiant d'abord les hypothèses : domaine, continuité, dérivabilité, puis calcul ou théorème utilisé.
\end{methodebox}

Sur $]0,+\infty[$, $x^\alpha=\e^{\alpha\ln x}$ est dérivable.


Le produit avec $\ln x$ est donc dérivable.


On a $(x^\alpha)'=\alpha x^{\alpha-1}$ et $(\ln x)'=1/x$.


Donc $f'(x)=\alpha x^{\alpha-1}\ln x+x^\alpha\frac1x$.


On factorise : $f'(x)=x^{\alpha-1}(\alpha\ln x+1)$.


Si $\alpha=0$, on retrouve $f(x)=\ln x$ et $f'(x)=1/x$, ce qui correspond à la formule.


\begin{attentionbox}
Vérification finale : la conclusion doit porter sur l'intervalle exact de travail et ne doit pas franchir un point exclu du domaine.
\end{attentionbox}

\end{corrbox}


\begin{corrbox}{9}{Arctangente composée}

\begin{methodebox}
On rédige en vérifiant d'abord les hypothèses : domaine, continuité, dérivabilité, puis calcul ou théorème utilisé.
\end{methodebox}

On a $u(x)=\frac{2x}{1-x^2}$.


La formule du quotient donne $u'(x)=\frac{2(1-x^2)+4x^2}{(1-x^2)^2}=\frac{2(1+x^2)}{(1-x^2)^2}$.


De plus $1+u^2=1+\frac{4x^2}{(1-x^2)^2}=\frac{(1-x^2)^2+4x^2}{(1-x^2)^2}$.


Le numérateur vaut $1+2x^2+x^4=(1+x^2)^2$.


Donc $1+u^2=\frac{(1+x^2)^2}{(1-x^2)^2}$.


Ainsi $f'(x)=\frac{u'}{1+u^2}=\frac{2}{1+x^2}$.


La dérivée simple obtenue rappelle les formules d'angle double pour la tangente.


\begin{attentionbox}
Vérification finale : la conclusion doit porter sur l'intervalle exact de travail et ne doit pas franchir un point exclu du domaine.
\end{attentionbox}

\end{corrbox}


\begin{corrbox}{10}{Fonctions hyperboliques}

\begin{methodebox}
On rédige en vérifiant d'abord les hypothèses : domaine, continuité, dérivabilité, puis calcul ou théorème utilisé.
\end{methodebox}

Pour tout $x$, $\ch x=\frac{\e^x+\e^{-x}}2>0$, donc $\ln(\ch x)$ est définie.


On a $(\ch x)'=\sh x$.


Donc $f'(x)=\frac{\sh x}{\ch x}=\thh x$.


Comme $\ch$ est paire, $\ln\circ\ch$ est paire.


Sur $[0,+\infty[$, $\sh x\ge0$ et $\ch x>0$, donc $f'(x)\ge0$.


La fonction est croissante sur $[0,+\infty[$.


\begin{attentionbox}
Vérification finale : la conclusion doit porter sur l'intervalle exact de travail et ne doit pas franchir un point exclu du domaine.
\end{attentionbox}

\end{corrbox}


\begin{corrbox}{11}{Dérivabilité d'une racine}

\begin{methodebox}
On rédige en vérifiant d'abord les hypothèses : domaine, continuité, dérivabilité, puis calcul ou théorème utilisé.
\end{methodebox}

On a $\sqrt{|x|}\to0$ lorsque $x\to0$, donc $f$ est continue en $0$.


Pour $h>0$, le quotient vaut $\frac{\sqrt h}{h}=\frac1{\sqrt h}$, qui tend vers $+\infty$.


Pour $h<0$, le quotient vaut $\frac{\sqrt{-h}}{h}=-\frac1{\sqrt{-h}}$, qui tend vers $-\infty$.


Il n'existe pas de limite finie du taux d'accroissement.


La fonction n'est donc pas dérivable en $0$.


Contrairement à $|x|$, le défaut n'est pas seulement un angle : la pente devient infinie.


\begin{attentionbox}
Vérification finale : la conclusion doit porter sur l'intervalle exact de travail et ne doit pas franchir un point exclu du domaine.
\end{attentionbox}

\end{corrbox}


\begin{corrbox}{12}{Dérivée d'une réciproque simple}

\begin{methodebox}
On rédige en vérifiant d'abord les hypothèses : domaine, continuité, dérivabilité, puis calcul ou théorème utilisé.
\end{methodebox}

La fonction est polynomiale donc continue sur $\R$.


Sa dérivée est $f'(x)=1+3x^2>0$.


Elle est donc strictement croissante.


De plus $f(x)\to-\infty$ quand $x\to-\infty$ et $f(x)\to+\infty$ quand $x\to+\infty$.


Par le théorème de la bijection continue, $f$ est une bijection de $\R$ sur $\R$.


Comme $f'$ ne s'annule jamais, la réciproque est dérivable et $(f^{-1})'(y)=\frac1{1+3(f^{-1}(y))^2}$.


Comme $f(0)=0$, on a $f^{-1}(0)=0$, donc $(f^{-1})'(0)=1$.


\begin{attentionbox}
Vérification finale : la conclusion doit porter sur l'intervalle exact de travail et ne doit pas franchir un point exclu du domaine.
\end{attentionbox}

\end{corrbox}


\begin{corrbox}{13}{Étude complète d'un quotient}

\begin{methodebox}
On rédige en vérifiant d'abord les hypothèses : domaine, continuité, dérivabilité, puis calcul ou théorème utilisé.
\end{methodebox}

La fonction est définie sur $\R\setminus\{-1,1\}$, sous réserve des restrictions évidentes imposées par les quotients, logarithmes ou réciproques trigonométriques.


Sur chaque intervalle ouvert de son domaine, elle est combinaison ou composée de fonctions usuelles dérivables.


Un calcul direct donne $f'(x)=\frac{-4x}{(x^2-1)^2}$.


La justification essentielle consiste à identifier la structure : quotient, composée, logarithme ou dérivation logarithmique.


Le signe de la dérivée se lit ensuite à partir des facteurs dont on connaît le signe.


On conclut les variations sur chaque intervalle du domaine, sans franchir les points exclus du domaine.


Cette dernière vérification évite l'erreur classique qui consiste à dresser un tableau de variations continu à travers une discontinuité.


\begin{attentionbox}
Vérification finale : la conclusion doit porter sur l'intervalle exact de travail et ne doit pas franchir un point exclu du domaine.
\end{attentionbox}

\end{corrbox}


\begin{corrbox}{14}{Logarithme et variations}

\begin{methodebox}
On rédige en vérifiant d'abord les hypothèses : domaine, continuité, dérivabilité, puis calcul ou théorème utilisé.
\end{methodebox}

La fonction est définie sur $]0,+\infty[$, sous réserve des restrictions évidentes imposées par les quotients, logarithmes ou réciproques trigonométriques.


Sur chaque intervalle ouvert de son domaine, elle est combinaison ou composée de fonctions usuelles dérivables.


Un calcul direct donne $f'(x)=\ln x$.


La justification essentielle consiste à identifier la structure : quotient, composée, logarithme ou dérivation logarithmique.


Le signe de la dérivée se lit ensuite à partir des facteurs dont on connaît le signe.


On conclut les variations sur chaque intervalle du domaine, sans franchir les points exclus du domaine.


Cette dernière vérification évite l'erreur classique qui consiste à dresser un tableau de variations continu à travers une discontinuité.


\begin{attentionbox}
Vérification finale : la conclusion doit porter sur l'intervalle exact de travail et ne doit pas franchir un point exclu du domaine.
\end{attentionbox}

\end{corrbox}


\begin{corrbox}{15}{Inégalité classique par étude de fonction}

\begin{methodebox}
On rédige en vérifiant d'abord les hypothèses : domaine, continuité, dérivabilité, puis calcul ou théorème utilisé.
\end{methodebox}

La fonction est définie sur $\R$, sous réserve des restrictions évidentes imposées par les quotients, logarithmes ou réciproques trigonométriques.


Sur chaque intervalle ouvert de son domaine, elle est combinaison ou composée de fonctions usuelles dérivables.


Un calcul direct donne $f'(x)=\e^x-1$.


La justification essentielle consiste à identifier la structure : quotient, composée, logarithme ou dérivation logarithmique.


Le signe de la dérivée se lit ensuite à partir des facteurs dont on connaît le signe.


On conclut les variations sur chaque intervalle du domaine, sans franchir les points exclus du domaine.


Cette dernière vérification évite l'erreur classique qui consiste à dresser un tableau de variations continu à travers une discontinuité.


\begin{attentionbox}
Vérification finale : la conclusion doit porter sur l'intervalle exact de travail et ne doit pas franchir un point exclu du domaine.
\end{attentionbox}

\end{corrbox}


\begin{corrbox}{16}{Puissance et logarithme}

\begin{methodebox}
On rédige en vérifiant d'abord les hypothèses : domaine, continuité, dérivabilité, puis calcul ou théorème utilisé.
\end{methodebox}

La fonction est définie sur $]0,+\infty[$, sous réserve des restrictions évidentes imposées par les quotients, logarithmes ou réciproques trigonométriques.


Sur chaque intervalle ouvert de son domaine, elle est combinaison ou composée de fonctions usuelles dérivables.


Un calcul direct donne $f'(x)=\frac{1-\ln x}{x^2}$.


La justification essentielle consiste à identifier la structure : quotient, composée, logarithme ou dérivation logarithmique.


Le signe de la dérivée se lit ensuite à partir des facteurs dont on connaît le signe.


On conclut les variations sur chaque intervalle du domaine, sans franchir les points exclus du domaine.


Cette dernière vérification évite l'erreur classique qui consiste à dresser un tableau de variations continu à travers une discontinuité.


\begin{attentionbox}
Vérification finale : la conclusion doit porter sur l'intervalle exact de travail et ne doit pas franchir un point exclu du domaine.
\end{attentionbox}

\end{corrbox}


\begin{corrbox}{17}{Tangente horizontale et extrema}

\begin{methodebox}
On rédige en vérifiant d'abord les hypothèses : domaine, continuité, dérivabilité, puis calcul ou théorème utilisé.
\end{methodebox}

La fonction est définie sur $\R$, sous réserve des restrictions évidentes imposées par les quotients, logarithmes ou réciproques trigonométriques.


Sur chaque intervalle ouvert de son domaine, elle est combinaison ou composée de fonctions usuelles dérivables.


Un calcul direct donne $f'(x)=3x^2-3$.


La justification essentielle consiste à identifier la structure : quotient, composée, logarithme ou dérivation logarithmique.


Le signe de la dérivée se lit ensuite à partir des facteurs dont on connaît le signe.


On conclut les variations sur chaque intervalle du domaine, sans franchir les points exclus du domaine.


Cette dernière vérification évite l'erreur classique qui consiste à dresser un tableau de variations continu à travers une discontinuité.


\begin{attentionbox}
Vérification finale : la conclusion doit porter sur l'intervalle exact de travail et ne doit pas franchir un point exclu du domaine.
\end{attentionbox}

\end{corrbox}


\begin{corrbox}{18}{Trigonométrie composée}

\begin{methodebox}
On rédige en vérifiant d'abord les hypothèses : domaine, continuité, dérivabilité, puis calcul ou théorème utilisé.
\end{methodebox}

La fonction est définie sur $\R$, sous réserve des restrictions évidentes imposées par les quotients, logarithmes ou réciproques trigonométriques.


Sur chaque intervalle ouvert de son domaine, elle est combinaison ou composée de fonctions usuelles dérivables.


Un calcul direct donne $f'(x)=\cos x-2\sin(2x)$.


La justification essentielle consiste à identifier la structure : quotient, composée, logarithme ou dérivation logarithmique.


Le signe de la dérivée se lit ensuite à partir des facteurs dont on connaît le signe.


On conclut les variations sur chaque intervalle du domaine, sans franchir les points exclus du domaine.


Cette dernière vérification évite l'erreur classique qui consiste à dresser un tableau de variations continu à travers une discontinuité.


\begin{attentionbox}
Vérification finale : la conclusion doit porter sur l'intervalle exact de travail et ne doit pas franchir un point exclu du domaine.
\end{attentionbox}

\end{corrbox}


\begin{corrbox}{19}{Arcsinus et domaine}

\begin{methodebox}
On rédige en vérifiant d'abord les hypothèses : domaine, continuité, dérivabilité, puis calcul ou théorème utilisé.
\end{methodebox}

La fonction est définie sur $[0,1]$, sous réserve des restrictions évidentes imposées par les quotients, logarithmes ou réciproques trigonométriques.


Sur chaque intervalle ouvert de son domaine, elle est combinaison ou composée de fonctions usuelles dérivables.


Un calcul direct donne $f'(x)=\frac{1}{\sqrt{x(1-x)}}\text{ sur }]0,1[$.


La justification essentielle consiste à identifier la structure : quotient, composée, logarithme ou dérivation logarithmique.


Le signe de la dérivée se lit ensuite à partir des facteurs dont on connaît le signe.


On conclut les variations sur chaque intervalle du domaine, sans franchir les points exclus du domaine.


Cette dernière vérification évite l'erreur classique qui consiste à dresser un tableau de variations continu à travers une discontinuité.


\begin{attentionbox}
Vérification finale : la conclusion doit porter sur l'intervalle exact de travail et ne doit pas franchir un point exclu du domaine.
\end{attentionbox}

\end{corrbox}


\begin{corrbox}{20}{Arccosinus et monotonie}

\begin{methodebox}
On rédige en vérifiant d'abord les hypothèses : domaine, continuité, dérivabilité, puis calcul ou théorème utilisé.
\end{methodebox}

La fonction est définie sur $\R$, sous réserve des restrictions évidentes imposées par les quotients, logarithmes ou réciproques trigonométriques.


Sur chaque intervalle ouvert de son domaine, elle est combinaison ou composée de fonctions usuelles dérivables.


Un calcul direct donne $f'(x)=-\frac{1}{1+x^2}$.


La justification essentielle consiste à identifier la structure : quotient, composée, logarithme ou dérivation logarithmique.


Le signe de la dérivée se lit ensuite à partir des facteurs dont on connaît le signe.


On conclut les variations sur chaque intervalle du domaine, sans franchir les points exclus du domaine.


Cette dernière vérification évite l'erreur classique qui consiste à dresser un tableau de variations continu à travers une discontinuité.


\begin{attentionbox}
Vérification finale : la conclusion doit porter sur l'intervalle exact de travail et ne doit pas franchir un point exclu du domaine.
\end{attentionbox}

\end{corrbox}


\begin{corrbox}{21}{Exponentielle rationnelle}

\begin{methodebox}
On rédige en vérifiant d'abord les hypothèses : domaine, continuité, dérivabilité, puis calcul ou théorème utilisé.
\end{methodebox}

La fonction est définie sur $\R$, sous réserve des restrictions évidentes imposées par les quotients, logarithmes ou réciproques trigonométriques.


Sur chaque intervalle ouvert de son domaine, elle est combinaison ou composée de fonctions usuelles dérivables.


Un calcul direct donne $f'(x)=\frac{\e^x}{(1+\e^x)^2}$.


La justification essentielle consiste à identifier la structure : quotient, composée, logarithme ou dérivation logarithmique.


Le signe de la dérivée se lit ensuite à partir des facteurs dont on connaît le signe.


On conclut les variations sur chaque intervalle du domaine, sans franchir les points exclus du domaine.


Cette dernière vérification évite l'erreur classique qui consiste à dresser un tableau de variations continu à travers une discontinuité.


\begin{attentionbox}
Vérification finale : la conclusion doit porter sur l'intervalle exact de travail et ne doit pas franchir un point exclu du domaine.
\end{attentionbox}

\end{corrbox}


\begin{corrbox}{22}{Fonction hyperbolique quotient}

\begin{methodebox}
On rédige en vérifiant d'abord les hypothèses : domaine, continuité, dérivabilité, puis calcul ou théorème utilisé.
\end{methodebox}

La fonction est définie sur $\R$, sous réserve des restrictions évidentes imposées par les quotients, logarithmes ou réciproques trigonométriques.


Sur chaque intervalle ouvert de son domaine, elle est combinaison ou composée de fonctions usuelles dérivables.


Un calcul direct donne $f'(x)=\frac{1}{\ch^2x}$.


La justification essentielle consiste à identifier la structure : quotient, composée, logarithme ou dérivation logarithmique.


Le signe de la dérivée se lit ensuite à partir des facteurs dont on connaît le signe.


On conclut les variations sur chaque intervalle du domaine, sans franchir les points exclus du domaine.


Cette dernière vérification évite l'erreur classique qui consiste à dresser un tableau de variations continu à travers une discontinuité.


\begin{attentionbox}
Vérification finale : la conclusion doit porter sur l'intervalle exact de travail et ne doit pas franchir un point exclu du domaine.
\end{attentionbox}

\end{corrbox}


\begin{corrbox}{23}{Dérivée logarithmique}

\begin{methodebox}
On rédige en vérifiant d'abord les hypothèses : domaine, continuité, dérivabilité, puis calcul ou théorème utilisé.
\end{methodebox}

La fonction est définie sur $\R$, sous réserve des restrictions évidentes imposées par les quotients, logarithmes ou réciproques trigonométriques.


Sur chaque intervalle ouvert de son domaine, elle est combinaison ou composée de fonctions usuelles dérivables.


Un calcul direct donne $f'(x)=(x^2+1)^x\left(\ln(x^2+1)+\frac{2x^2}{x^2+1}\right)$.


La justification essentielle consiste à identifier la structure : quotient, composée, logarithme ou dérivation logarithmique.


Le signe de la dérivée se lit ensuite à partir des facteurs dont on connaît le signe.


On conclut les variations sur chaque intervalle du domaine, sans franchir les points exclus du domaine.


Cette dernière vérification évite l'erreur classique qui consiste à dresser un tableau de variations continu à travers une discontinuité.


\begin{attentionbox}
Vérification finale : la conclusion doit porter sur l'intervalle exact de travail et ne doit pas franchir un point exclu du domaine.
\end{attentionbox}

\end{corrbox}


\begin{corrbox}{24}{Optimisation élémentaire}

\begin{methodebox}
On rédige en vérifiant d'abord les hypothèses : domaine, continuité, dérivabilité, puis calcul ou théorème utilisé.
\end{methodebox}

La fonction est définie sur $]0,1[$, sous réserve des restrictions évidentes imposées par les quotients, logarithmes ou réciproques trigonométriques.


Sur chaque intervalle ouvert de son domaine, elle est combinaison ou composée de fonctions usuelles dérivables.


Un calcul direct donne $f'(x)=x^a(1-x)^b\left(\frac{a}{x}-\frac{b}{1-x}\right)$.


La justification essentielle consiste à identifier la structure : quotient, composée, logarithme ou dérivation logarithmique.


Le signe de la dérivée se lit ensuite à partir des facteurs dont on connaît le signe.


On conclut les variations sur chaque intervalle du domaine, sans franchir les points exclus du domaine.


Cette dernière vérification évite l'erreur classique qui consiste à dresser un tableau de variations continu à travers une discontinuité.


\begin{attentionbox}
Vérification finale : la conclusion doit porter sur l'intervalle exact de travail et ne doit pas franchir un point exclu du domaine.
\end{attentionbox}

\end{corrbox}


\begin{corrbox}{25}{Dérivabilité implique continuité}

\begin{methodebox}
On rédige en vérifiant d'abord les hypothèses : domaine, continuité, dérivabilité, puis calcul ou théorème utilisé.
\end{methodebox}

Comme $f$ est dérivable en $a$, le quotient $\frac{f(a+h)-f(a)}h$ tend vers $f'(a)$.


Pour $h
eq0$, on pose $\varepsilon(h)=\frac{f(a+h)-f(a)}h-f'(a)$.


Alors $\varepsilon(h)\to0$.


On obtient $f(a+h)-f(a)=h f'(a)+h\varepsilon(h)$.


Le membre de droite tend vers $0$, donc $f(a+h)\to f(a)$.


La fonction $x\mapsto |x|$ est continue en $0$ mais non dérivable en $0$, ce qui réfute la réciproque.


\begin{attentionbox}
Vérification finale : la conclusion doit porter sur l'intervalle exact de travail et ne doit pas franchir un point exclu du domaine.
\end{attentionbox}

\end{corrbox}


\begin{corrbox}{26}{Formule du produit par les taux}

\begin{methodebox}
On rédige en vérifiant d'abord les hypothèses : domaine, continuité, dérivabilité, puis calcul ou théorème utilisé.
\end{methodebox}

On écrit $f(a+h)g(a+h)-f(a)g(a)$ comme $(f(a+h)-f(a))g(a+h)+f(a)(g(a+h)-g(a))$.


Après division par $h$, on obtient deux termes.


Le premier tend vers $f'(a)g(a)$ car $g$ est continue en $a$.


Le second tend vers $f(a)g'(a)$.


La somme des limites donne la formule annoncée.


La continuité de $g$ n'est pas une hypothèse supplémentaire : elle découle de sa dérivabilité.


\begin{attentionbox}
Vérification finale : la conclusion doit porter sur l'intervalle exact de travail et ne doit pas franchir un point exclu du domaine.
\end{attentionbox}

\end{corrbox}


\begin{corrbox}{27}{Formule de la composée}

\begin{methodebox}
On rédige en vérifiant d'abord les hypothèses : domaine, continuité, dérivabilité, puis calcul ou théorème utilisé.
\end{methodebox}

La dérivabilité de $f$ en $g(a)$ donne $f(g(a)+k)=f(g(a))+kf'(g(a))+k\varepsilon(k)$ avec $\varepsilon(k)\to0$.


On pose $k=g(a+h)-g(a)$.


Comme $g$ est dérivable, elle est continue, donc $k\to0$.


On divise l'égalité obtenue par $h$.


Le quotient $\frac{g(a+h)-g(a)}h$ tend vers $g'(a)$.


Le terme avec $\varepsilon(k)$ tend vers $0$ car il est produit d'un quotient borné au voisinage de la limite et d'un terme tendant vers $0$.


La formule de la composée est démontrée.


\begin{attentionbox}
Vérification finale : la conclusion doit porter sur l'intervalle exact de travail et ne doit pas franchir un point exclu du domaine.
\end{attentionbox}

\end{corrbox}


\begin{corrbox}{28}{Dérivée de la réciproque}

\begin{methodebox}
On rédige en vérifiant d'abord les hypothèses : domaine, continuité, dérivabilité, puis calcul ou théorème utilisé.
\end{methodebox}

On considère $\frac{f^{-1}(y)-f^{-1}(b)}{y-b}$.


En posant $x=f^{-1}(y)$, on a $y=f(x)$ et $f^{-1}(b)=a$.


Le quotient devient $\frac{x-a}{f(x)-f(a)}$.


Lorsque $y\to b$, la continuité de $f^{-1}$ donne $x\to a$.


Le quotient est donc l'inverse de $\frac{f(x)-f(a)}{x-a}$.


Comme ce dernier tend vers $f'(a)
e0$, la limite vaut $1/f'(a)$.


Ainsi $(f^{-1})'(b)=\frac1{f'(a)}$.


\begin{attentionbox}
Vérification finale : la conclusion doit porter sur l'intervalle exact de travail et ne doit pas franchir un point exclu du domaine.
\end{attentionbox}

\end{corrbox}


\begin{corrbox}{29}{Condition nécessaire d'extremum}

\begin{methodebox}
On rédige en vérifiant d'abord les hypothèses : domaine, continuité, dérivabilité, puis calcul ou théorème utilisé.
\end{methodebox}

Pour un maximum local, on a $f(a+h)\le f(a)$ pour $h$ assez petit.


Si $h>0$, le quotient $\frac{f(a+h)-f(a)}h$ est négatif ou nul.


En passant à la limite à droite, on obtient $f'_d(a)\le0$.


Si $h<0$, le numérateur est encore négatif ou nul, mais la division par $h$ inverse le sens : le quotient est positif ou nul.


La limite à gauche donne $f'_g(a)\ge0$.


La dérivabilité impose l'égalité des deux dérivées latérales, donc $f'(a)=0$.


Pour un minimum local, les inégalités sont inversées et la conclusion reste identique.


\begin{attentionbox}
Vérification finale : la conclusion doit porter sur l'intervalle exact de travail et ne doit pas franchir un point exclu du domaine.
\end{attentionbox}

\end{corrbox}


\begin{corrbox}{30}{Théorème de Rolle}

\begin{methodebox}
On rédige en vérifiant d'abord les hypothèses : domaine, continuité, dérivabilité, puis calcul ou théorème utilisé.
\end{methodebox}

Soit $f$ continue sur $[a,b]$, dérivable sur $]a,b[$ et telle que $f(a)=f(b)$.


Par continuité sur un segment, $f$ atteint un minimum et un maximum.


Si ces deux valeurs sont égales, $f$ est constante et n'importe quel point intérieur convient.


Sinon, au moins l'un des extrema n'est pas égal à la valeur commune $f(a)=f(b)$.


Cet extremum ne peut donc pas être atteint seulement aux deux extrémités.


Il existe donc un point intérieur $c$ où $f$ admet un extremum local.


Comme $f$ est dérivable en $c$, la condition nécessaire donne $f'(c)=0$.


\begin{attentionbox}
Vérification finale : la conclusion doit porter sur l'intervalle exact de travail et ne doit pas franchir un point exclu du domaine.
\end{attentionbox}

\end{corrbox}


\begin{corrbox}{31}{Théorème des accroissements finis}

\begin{methodebox}
On rédige en vérifiant d'abord les hypothèses : domaine, continuité, dérivabilité, puis calcul ou théorème utilisé.
\end{methodebox}

On définit $g(x)=f(x)-\frac{f(b)-f(a)}{b-a}(x-a)$.


La fonction $g$ est continue sur $[a,b]$ et dérivable sur $]a,b[$.


On a $g(a)=f(a)$ et $g(b)=f(b)-(f(b)-f(a))=f(a)$.


Par Rolle, il existe $c\in]a,b[$ tel que $g'(c)=0$.


Or $g'(x)=f'(x)-\frac{f(b)-f(a)}{b-a}$.


Donc $f'(c)=\frac{f(b)-f(a)}{b-a}$.


\begin{attentionbox}
Vérification finale : la conclusion doit porter sur l'intervalle exact de travail et ne doit pas franchir un point exclu du domaine.
\end{attentionbox}

\end{corrbox}


\begin{corrbox}{32}{Inégalité des accroissements finis}

\begin{methodebox}
On rédige en vérifiant d'abord les hypothèses : domaine, continuité, dérivabilité, puis calcul ou théorème utilisé.
\end{methodebox}

Pour $x<y$, le segment $[x,y]$ est inclus dans $I$ car $I$ est un intervalle.


Par le TAF, il existe $c\in]x,y[$ tel que $f(y)-f(x)=f'(c)(y-x)$.


En valeur absolue, $|f(y)-f(x)|=|f'(c)|\,|y-x|$.


La majoration $|f'(c)|\le K$ donne $|f(y)-f(x)|\le K|y-x|$.


Le cas $x=y$ est immédiat, et le cas $y<x$ se déduit par symétrie.


\begin{attentionbox}
Vérification finale : la conclusion doit porter sur l'intervalle exact de travail et ne doit pas franchir un point exclu du domaine.
\end{attentionbox}

\end{corrbox}


\begin{corrbox}{33}{Rolle et racines successives}

\begin{methodebox}
On rédige en vérifiant d'abord les hypothèses : domaine, continuité, dérivabilité, puis calcul ou théorème utilisé.
\end{methodebox}

Notons $a_1<\cdots<a_n$ les racines distinctes de $P$.


Sur chaque segment $[a_i,a_{i+1}]$, le polynôme est continu et dérivable.


On a $P(a_i)=P(a_{i+1})=0$.


Par Rolle, il existe $c_i\in]a_i,a_{i+1}[$ tel que $P'(c_i)=0$.


Les intervalles étant disjoints, les $c_i$ sont distincts.


Ainsi $P'$ possède au moins $n-1$ racines réelles distinctes.


Les racines de la dérivée s'intercalent entre les racines de $P$.


\begin{attentionbox}
Vérification finale : la conclusion doit porter sur l'intervalle exact de travail et ne doit pas franchir un point exclu du domaine.
\end{attentionbox}

\end{corrbox}


\begin{corrbox}{34}{Convexité sans nom par dérivée croissante}

\begin{methodebox}
On rédige en vérifiant d'abord les hypothèses : domaine, continuité, dérivabilité, puis calcul ou théorème utilisé.
\end{methodebox}

Par le TAF sur $[a,x]$, il existe $u\in]a,x[$ tel que $\frac{f(x)-f(a)}{x-a}=f'(u)$.


Par le TAF sur $[x,b]$, il existe $v\in]x,b[$ tel que $\frac{f(b)-f(x)}{b-x}=f'(v)$.


Comme $u<v$ et $f'$ est croissante, on a $f'(u)\le f'(v)$.


Donc $\frac{f(x)-f(a)}{x-a}\le \frac{f(b)-f(x)}{b-x}$.


En multipliant par les quantités positives, on obtient $(b-x)(f(x)-f(a))\le(x-a)(f(b)-f(x))$.


Après réarrangement, $f(x)\le \frac{b-x}{b-a}f(a)+\frac{x-a}{b-a}f(b)$.


\begin{attentionbox}
Vérification finale : la conclusion doit porter sur l'intervalle exact de travail et ne doit pas franchir un point exclu du domaine.
\end{attentionbox}

\end{corrbox}


\begin{corrbox}{35}{Inégalité $\ln x\le x-1$ et raffinements}

\begin{methodebox}
On rédige en vérifiant d'abord les hypothèses : domaine, continuité, dérivabilité, puis calcul ou théorème utilisé.
\end{methodebox}

On commence par identifier le domaine naturel de la fonction : logarithme strictement positif, quotient non nul, ou domaine des fonctions réciproques trigonométriques.


Sur chaque intervalle admissible, les fonctions usuelles impliquées sont dérivables.


Le calcul de la dérivée s'effectue par les règles de somme, produit, quotient ou composée.


Le signe de la dérivée permet d'obtenir les variations par le théorème des accroissements finis.


Un extremum global obtenu par le tableau de variations fournit l'inégalité demandée.


Lorsqu'il s'agit d'une équation, la continuité donne l'existence par changement de signe ou limites aux bornes.


La stricte monotonie donne l'unicité, car une fonction strictement monotone est injective.


La conclusion est donc entièrement justifiée par les théorèmes du chapitre, sans argument graphique.


\begin{attentionbox}
Vérification finale : la conclusion doit porter sur l'intervalle exact de travail et ne doit pas franchir un point exclu du domaine.
\end{attentionbox}

\end{corrbox}


\begin{corrbox}{36}{Encadrement de l'exponentielle}

\begin{methodebox}
On rédige en vérifiant d'abord les hypothèses : domaine, continuité, dérivabilité, puis calcul ou théorème utilisé.
\end{methodebox}

On commence par identifier le domaine naturel de la fonction : logarithme strictement positif, quotient non nul, ou domaine des fonctions réciproques trigonométriques.


Sur chaque intervalle admissible, les fonctions usuelles impliquées sont dérivables.


Le calcul de la dérivée s'effectue par les règles de somme, produit, quotient ou composée.


Le signe de la dérivée permet d'obtenir les variations par le théorème des accroissements finis.


Un extremum global obtenu par le tableau de variations fournit l'inégalité demandée.


Lorsqu'il s'agit d'une équation, la continuité donne l'existence par changement de signe ou limites aux bornes.


La stricte monotonie donne l'unicité, car une fonction strictement monotone est injective.


La conclusion est donc entièrement justifiée par les théorèmes du chapitre, sans argument graphique.


\begin{attentionbox}
Vérification finale : la conclusion doit porter sur l'intervalle exact de travail et ne doit pas franchir un point exclu du domaine.
\end{attentionbox}

\end{corrbox}


\begin{corrbox}{37}{Unicité pour une équation transcendante}

\begin{methodebox}
On rédige en vérifiant d'abord les hypothèses : domaine, continuité, dérivabilité, puis calcul ou théorème utilisé.
\end{methodebox}

On commence par identifier le domaine naturel de la fonction : logarithme strictement positif, quotient non nul, ou domaine des fonctions réciproques trigonométriques.


Sur chaque intervalle admissible, les fonctions usuelles impliquées sont dérivables.


Le calcul de la dérivée s'effectue par les règles de somme, produit, quotient ou composée.


Le signe de la dérivée permet d'obtenir les variations par le théorème des accroissements finis.


Un extremum global obtenu par le tableau de variations fournit l'inégalité demandée.


Lorsqu'il s'agit d'une équation, la continuité donne l'existence par changement de signe ou limites aux bornes.


La stricte monotonie donne l'unicité, car une fonction strictement monotone est injective.


La conclusion est donc entièrement justifiée par les théorèmes du chapitre, sans argument graphique.


\begin{attentionbox}
Vérification finale : la conclusion doit porter sur l'intervalle exact de travail et ne doit pas franchir un point exclu du domaine.
\end{attentionbox}

\end{corrbox}


\begin{corrbox}{38}{Point fixe de $\cos$}

\begin{methodebox}
On rédige en vérifiant d'abord les hypothèses : domaine, continuité, dérivabilité, puis calcul ou théorème utilisé.
\end{methodebox}

On commence par identifier le domaine naturel de la fonction : logarithme strictement positif, quotient non nul, ou domaine des fonctions réciproques trigonométriques.


Sur chaque intervalle admissible, les fonctions usuelles impliquées sont dérivables.


Le calcul de la dérivée s'effectue par les règles de somme, produit, quotient ou composée.


Le signe de la dérivée permet d'obtenir les variations par le théorème des accroissements finis.


Un extremum global obtenu par le tableau de variations fournit l'inégalité demandée.


Lorsqu'il s'agit d'une équation, la continuité donne l'existence par changement de signe ou limites aux bornes.


La stricte monotonie donne l'unicité, car une fonction strictement monotone est injective.


La conclusion est donc entièrement justifiée par les théorèmes du chapitre, sans argument graphique.


\begin{attentionbox}
Vérification finale : la conclusion doit porter sur l'intervalle exact de travail et ne doit pas franchir un point exclu du domaine.
\end{attentionbox}

\end{corrbox}


\begin{corrbox}{39}{Étude de $x-\arctan x$}

\begin{methodebox}
On rédige en vérifiant d'abord les hypothèses : domaine, continuité, dérivabilité, puis calcul ou théorème utilisé.
\end{methodebox}

On commence par identifier le domaine naturel de la fonction : logarithme strictement positif, quotient non nul, ou domaine des fonctions réciproques trigonométriques.


Sur chaque intervalle admissible, les fonctions usuelles impliquées sont dérivables.


Le calcul de la dérivée s'effectue par les règles de somme, produit, quotient ou composée.


Le signe de la dérivée permet d'obtenir les variations par le théorème des accroissements finis.


Un extremum global obtenu par le tableau de variations fournit l'inégalité demandée.


Lorsqu'il s'agit d'une équation, la continuité donne l'existence par changement de signe ou limites aux bornes.


La stricte monotonie donne l'unicité, car une fonction strictement monotone est injective.


La conclusion est donc entièrement justifiée par les théorèmes du chapitre, sans argument graphique.


\begin{attentionbox}
Vérification finale : la conclusion doit porter sur l'intervalle exact de travail et ne doit pas franchir un point exclu du domaine.
\end{attentionbox}

\end{corrbox}


\begin{corrbox}{40}{Étude de $\sin x/x$}

\begin{methodebox}
On rédige en vérifiant d'abord les hypothèses : domaine, continuité, dérivabilité, puis calcul ou théorème utilisé.
\end{methodebox}

On commence par identifier le domaine naturel de la fonction : logarithme strictement positif, quotient non nul, ou domaine des fonctions réciproques trigonométriques.


Sur chaque intervalle admissible, les fonctions usuelles impliquées sont dérivables.


Le calcul de la dérivée s'effectue par les règles de somme, produit, quotient ou composée.


Le signe de la dérivée permet d'obtenir les variations par le théorème des accroissements finis.


Un extremum global obtenu par le tableau de variations fournit l'inégalité demandée.


Lorsqu'il s'agit d'une équation, la continuité donne l'existence par changement de signe ou limites aux bornes.


La stricte monotonie donne l'unicité, car une fonction strictement monotone est injective.


La conclusion est donc entièrement justifiée par les théorèmes du chapitre, sans argument graphique.


\begin{attentionbox}
Vérification finale : la conclusion doit porter sur l'intervalle exact de travail et ne doit pas franchir un point exclu du domaine.
\end{attentionbox}

\end{corrbox}


\begin{corrbox}{41}{Maximum d'une fonction avec paramètre}

\begin{methodebox}
On rédige en vérifiant d'abord les hypothèses : domaine, continuité, dérivabilité, puis calcul ou théorème utilisé.
\end{methodebox}

On commence par identifier le domaine naturel de la fonction : logarithme strictement positif, quotient non nul, ou domaine des fonctions réciproques trigonométriques.


Sur chaque intervalle admissible, les fonctions usuelles impliquées sont dérivables.


Le calcul de la dérivée s'effectue par les règles de somme, produit, quotient ou composée.


Le signe de la dérivée permet d'obtenir les variations par le théorème des accroissements finis.


Un extremum global obtenu par le tableau de variations fournit l'inégalité demandée.


Lorsqu'il s'agit d'une équation, la continuité donne l'existence par changement de signe ou limites aux bornes.


La stricte monotonie donne l'unicité, car une fonction strictement monotone est injective.


La conclusion est donc entièrement justifiée par les théorèmes du chapitre, sans argument graphique.


\begin{attentionbox}
Vérification finale : la conclusion doit porter sur l'intervalle exact de travail et ne doit pas franchir un point exclu du domaine.
\end{attentionbox}

\end{corrbox}


\begin{corrbox}{42}{Réciproque de $x+\e^x$}

\begin{methodebox}
On rédige en vérifiant d'abord les hypothèses : domaine, continuité, dérivabilité, puis calcul ou théorème utilisé.
\end{methodebox}

On commence par identifier le domaine naturel de la fonction : logarithme strictement positif, quotient non nul, ou domaine des fonctions réciproques trigonométriques.


Sur chaque intervalle admissible, les fonctions usuelles impliquées sont dérivables.


Le calcul de la dérivée s'effectue par les règles de somme, produit, quotient ou composée.


Le signe de la dérivée permet d'obtenir les variations par le théorème des accroissements finis.


Un extremum global obtenu par le tableau de variations fournit l'inégalité demandée.


Lorsqu'il s'agit d'une équation, la continuité donne l'existence par changement de signe ou limites aux bornes.


La stricte monotonie donne l'unicité, car une fonction strictement monotone est injective.


La conclusion est donc entièrement justifiée par les théorèmes du chapitre, sans argument graphique.


\begin{attentionbox}
Vérification finale : la conclusion doit porter sur l'intervalle exact de travail et ne doit pas franchir un point exclu du domaine.
\end{attentionbox}

\end{corrbox}


\begin{corrbox}{43}{Comparaison logarithme-puissance}

\begin{methodebox}
On rédige en vérifiant d'abord les hypothèses : domaine, continuité, dérivabilité, puis calcul ou théorème utilisé.
\end{methodebox}

On commence par identifier le domaine naturel de la fonction : logarithme strictement positif, quotient non nul, ou domaine des fonctions réciproques trigonométriques.


Sur chaque intervalle admissible, les fonctions usuelles impliquées sont dérivables.


Le calcul de la dérivée s'effectue par les règles de somme, produit, quotient ou composée.


Le signe de la dérivée permet d'obtenir les variations par le théorème des accroissements finis.


Un extremum global obtenu par le tableau de variations fournit l'inégalité demandée.


Lorsqu'il s'agit d'une équation, la continuité donne l'existence par changement de signe ou limites aux bornes.


La stricte monotonie donne l'unicité, car une fonction strictement monotone est injective.


La conclusion est donc entièrement justifiée par les théorèmes du chapitre, sans argument graphique.


\begin{attentionbox}
Vérification finale : la conclusion doit porter sur l'intervalle exact de travail et ne doit pas franchir un point exclu du domaine.
\end{attentionbox}

\end{corrbox}


\begin{corrbox}{44}{Hyperboliques et inégalités}

\begin{methodebox}
On rédige en vérifiant d'abord les hypothèses : domaine, continuité, dérivabilité, puis calcul ou théorème utilisé.
\end{methodebox}

On commence par identifier le domaine naturel de la fonction : logarithme strictement positif, quotient non nul, ou domaine des fonctions réciproques trigonométriques.


Sur chaque intervalle admissible, les fonctions usuelles impliquées sont dérivables.


Le calcul de la dérivée s'effectue par les règles de somme, produit, quotient ou composée.


Le signe de la dérivée permet d'obtenir les variations par le théorème des accroissements finis.


Un extremum global obtenu par le tableau de variations fournit l'inégalité demandée.


Lorsqu'il s'agit d'une équation, la continuité donne l'existence par changement de signe ou limites aux bornes.


La stricte monotonie donne l'unicité, car une fonction strictement monotone est injective.


La conclusion est donc entièrement justifiée par les théorèmes du chapitre, sans argument graphique.


\begin{attentionbox}
Vérification finale : la conclusion doit porter sur l'intervalle exact de travail et ne doit pas franchir un point exclu du domaine.
\end{attentionbox}

\end{corrbox}


\begin{corrbox}{45}{Problème CCP -- Rolle itéré et polynômes}

\begin{methodebox}
On rédige en vérifiant d'abord les hypothèses : domaine, continuité, dérivabilité, puis calcul ou théorème utilisé.
\end{methodebox}

La première étape consiste à isoler les intervalles où tous les objets sont bien définis.


La continuité sur les segments et la dérivabilité sur les ouverts sont les hypothèses minimales permettant d'utiliser Rolle ou le TAF.


Une fois ces hypothèses vérifiées, on calcule la dérivée et on cherche une factorisation exploitable.


Dans un problème de concours, cette factorisation est souvent l'idée centrale : elle transforme un calcul long en signe lisible.


Le théorème de Rolle sert à faire apparaître un point où la dérivée s'annule entre deux valeurs égales.


Le théorème des accroissements finis relie une variation globale à une dérivée locale.


L'inégalité des accroissements finis transforme une borne sur la dérivée en contrôle lipschitzien.


La dernière question se conclut par une phrase complète : existence par continuité, unicité par monotonie, ou borne par extremum global.


Toutes les hypothèses du théorème utilisé ont été vérifiées : c'est ce qui donne une rédaction de niveau concours.


\begin{attentionbox}
Vérification finale : la conclusion doit porter sur l'intervalle exact de travail et ne doit pas franchir un point exclu du domaine.
\end{attentionbox}

\end{corrbox}


\begin{corrbox}{46}{Problème Mines -- Inégalités par TAF}

\begin{methodebox}
On rédige en vérifiant d'abord les hypothèses : domaine, continuité, dérivabilité, puis calcul ou théorème utilisé.
\end{methodebox}

La première étape consiste à isoler les intervalles où tous les objets sont bien définis.


La continuité sur les segments et la dérivabilité sur les ouverts sont les hypothèses minimales permettant d'utiliser Rolle ou le TAF.


Une fois ces hypothèses vérifiées, on calcule la dérivée et on cherche une factorisation exploitable.


Dans un problème de concours, cette factorisation est souvent l'idée centrale : elle transforme un calcul long en signe lisible.


Le théorème de Rolle sert à faire apparaître un point où la dérivée s'annule entre deux valeurs égales.


Le théorème des accroissements finis relie une variation globale à une dérivée locale.


L'inégalité des accroissements finis transforme une borne sur la dérivée en contrôle lipschitzien.


La dernière question se conclut par une phrase complète : existence par continuité, unicité par monotonie, ou borne par extremum global.


Toutes les hypothèses du théorème utilisé ont été vérifiées : c'est ce qui donne une rédaction de niveau concours.


\begin{attentionbox}
Vérification finale : la conclusion doit porter sur l'intervalle exact de travail et ne doit pas franchir un point exclu du domaine.
\end{attentionbox}

\end{corrbox}


\begin{corrbox}{47}{Problème Centrale -- Réciproque et paramètre}

\begin{methodebox}
On rédige en vérifiant d'abord les hypothèses : domaine, continuité, dérivabilité, puis calcul ou théorème utilisé.
\end{methodebox}

La première étape consiste à isoler les intervalles où tous les objets sont bien définis.


La continuité sur les segments et la dérivabilité sur les ouverts sont les hypothèses minimales permettant d'utiliser Rolle ou le TAF.


Une fois ces hypothèses vérifiées, on calcule la dérivée et on cherche une factorisation exploitable.


Dans un problème de concours, cette factorisation est souvent l'idée centrale : elle transforme un calcul long en signe lisible.


Le théorème de Rolle sert à faire apparaître un point où la dérivée s'annule entre deux valeurs égales.


Le théorème des accroissements finis relie une variation globale à une dérivée locale.


L'inégalité des accroissements finis transforme une borne sur la dérivée en contrôle lipschitzien.


La dernière question se conclut par une phrase complète : existence par continuité, unicité par monotonie, ou borne par extremum global.


Toutes les hypothèses du théorème utilisé ont été vérifiées : c'est ce qui donne une rédaction de niveau concours.


\begin{attentionbox}
Vérification finale : la conclusion doit porter sur l'intervalle exact de travail et ne doit pas franchir un point exclu du domaine.
\end{attentionbox}

\end{corrbox}


\begin{corrbox}{48}{Problème oral -- Tangentes communes}

\begin{methodebox}
On rédige en vérifiant d'abord les hypothèses : domaine, continuité, dérivabilité, puis calcul ou théorème utilisé.
\end{methodebox}

La première étape consiste à isoler les intervalles où tous les objets sont bien définis.


La continuité sur les segments et la dérivabilité sur les ouverts sont les hypothèses minimales permettant d'utiliser Rolle ou le TAF.


Une fois ces hypothèses vérifiées, on calcule la dérivée et on cherche une factorisation exploitable.


Dans un problème de concours, cette factorisation est souvent l'idée centrale : elle transforme un calcul long en signe lisible.


Le théorème de Rolle sert à faire apparaître un point où la dérivée s'annule entre deux valeurs égales.


Le théorème des accroissements finis relie une variation globale à une dérivée locale.


L'inégalité des accroissements finis transforme une borne sur la dérivée en contrôle lipschitzien.


La dernière question se conclut par une phrase complète : existence par continuité, unicité par monotonie, ou borne par extremum global.


Toutes les hypothèses du théorème utilisé ont été vérifiées : c'est ce qui donne une rédaction de niveau concours.


\begin{attentionbox}
Vérification finale : la conclusion doit porter sur l'intervalle exact de travail et ne doit pas franchir un point exclu du domaine.
\end{attentionbox}

\end{corrbox}


\begin{corrbox}{49}{Problème oral -- Arctangente et logarithme}

\begin{methodebox}
On rédige en vérifiant d'abord les hypothèses : domaine, continuité, dérivabilité, puis calcul ou théorème utilisé.
\end{methodebox}

La première étape consiste à isoler les intervalles où tous les objets sont bien définis.


La continuité sur les segments et la dérivabilité sur les ouverts sont les hypothèses minimales permettant d'utiliser Rolle ou le TAF.


Une fois ces hypothèses vérifiées, on calcule la dérivée et on cherche une factorisation exploitable.


Dans un problème de concours, cette factorisation est souvent l'idée centrale : elle transforme un calcul long en signe lisible.


Le théorème de Rolle sert à faire apparaître un point où la dérivée s'annule entre deux valeurs égales.


Le théorème des accroissements finis relie une variation globale à une dérivée locale.


L'inégalité des accroissements finis transforme une borne sur la dérivée en contrôle lipschitzien.


La dernière question se conclut par une phrase complète : existence par continuité, unicité par monotonie, ou borne par extremum global.


Toutes les hypothèses du théorème utilisé ont été vérifiées : c'est ce qui donne une rédaction de niveau concours.


\begin{attentionbox}
Vérification finale : la conclusion doit porter sur l'intervalle exact de travail et ne doit pas franchir un point exclu du domaine.
\end{attentionbox}

\end{corrbox}


\begin{corrbox}{50}{Problème bilan -- Fonction définie par raccords}

\begin{methodebox}
On rédige en vérifiant d'abord les hypothèses : domaine, continuité, dérivabilité, puis calcul ou théorème utilisé.
\end{methodebox}

La première étape consiste à isoler les intervalles où tous les objets sont bien définis.


La continuité sur les segments et la dérivabilité sur les ouverts sont les hypothèses minimales permettant d'utiliser Rolle ou le TAF.


Une fois ces hypothèses vérifiées, on calcule la dérivée et on cherche une factorisation exploitable.


Dans un problème de concours, cette factorisation est souvent l'idée centrale : elle transforme un calcul long en signe lisible.


Le théorème de Rolle sert à faire apparaître un point où la dérivée s'annule entre deux valeurs égales.


Le théorème des accroissements finis relie une variation globale à une dérivée locale.


L'inégalité des accroissements finis transforme une borne sur la dérivée en contrôle lipschitzien.


La dernière question se conclut par une phrase complète : existence par continuité, unicité par monotonie, ou borne par extremum global.


Toutes les hypothèses du théorème utilisé ont été vérifiées : c'est ce qui donne une rédaction de niveau concours.


\begin{attentionbox}
Vérification finale : la conclusion doit porter sur l'intervalle exact de travail et ne doit pas franchir un point exclu du domaine.
\end{attentionbox}

\end{corrbox}


\newpage
\section{Annexe -- Grilles de rédaction pour copies de concours}
\subsection*{Grille 1 -- réflexe de rédaction}

\begin{methodebox}

Avant de dériver, écrire le domaine de définition.

Avant d'appliquer Rolle, vérifier la continuité sur le segment.

Avant d'appliquer Rolle, vérifier la dérivabilité sur l'intervalle ouvert.

Avant d'appliquer le TAF, vérifier que les deux bornes appartiennent au même intervalle.

Avant d'utiliser une réciproque, vérifier la bijectivité par continuité et stricte monotonie.

Avant d'utiliser la formule de la réciproque, vérifier que la dérivée ne s'annule pas.

Pour une inégalité, construire une fonction auxiliaire dont on étudie le minimum.

Pour une unicité, chercher une stricte monotonie ou une injectivité.

Pour une existence, chercher un changement de signe ou une limite aux bornes.

Pour une borne, chercher une majoration uniforme de la dérivée.

\end{methodebox}

\subsection*{Grille 2 -- réflexe de rédaction}

\begin{methodebox}

Avant de dériver, écrire le domaine de définition.

Avant d'appliquer Rolle, vérifier la continuité sur le segment.

Avant d'appliquer Rolle, vérifier la dérivabilité sur l'intervalle ouvert.

Avant d'appliquer le TAF, vérifier que les deux bornes appartiennent au même intervalle.

Avant d'utiliser une réciproque, vérifier la bijectivité par continuité et stricte monotonie.

Avant d'utiliser la formule de la réciproque, vérifier que la dérivée ne s'annule pas.

Pour une inégalité, construire une fonction auxiliaire dont on étudie le minimum.

Pour une unicité, chercher une stricte monotonie ou une injectivité.

Pour une existence, chercher un changement de signe ou une limite aux bornes.

Pour une borne, chercher une majoration uniforme de la dérivée.

\end{methodebox}

\subsection*{Grille 3 -- réflexe de rédaction}

\begin{methodebox}

Avant de dériver, écrire le domaine de définition.

Avant d'appliquer Rolle, vérifier la continuité sur le segment.

Avant d'appliquer Rolle, vérifier la dérivabilité sur l'intervalle ouvert.

Avant d'appliquer le TAF, vérifier que les deux bornes appartiennent au même intervalle.

Avant d'utiliser une réciproque, vérifier la bijectivité par continuité et stricte monotonie.

Avant d'utiliser la formule de la réciproque, vérifier que la dérivée ne s'annule pas.

Pour une inégalité, construire une fonction auxiliaire dont on étudie le minimum.

Pour une unicité, chercher une stricte monotonie ou une injectivité.

Pour une existence, chercher un changement de signe ou une limite aux bornes.

Pour une borne, chercher une majoration uniforme de la dérivée.

\end{methodebox}

\subsection*{Grille 4 -- réflexe de rédaction}

\begin{methodebox}

Avant de dériver, écrire le domaine de définition.

Avant d'appliquer Rolle, vérifier la continuité sur le segment.

Avant d'appliquer Rolle, vérifier la dérivabilité sur l'intervalle ouvert.

Avant d'appliquer le TAF, vérifier que les deux bornes appartiennent au même intervalle.

Avant d'utiliser une réciproque, vérifier la bijectivité par continuité et stricte monotonie.

Avant d'utiliser la formule de la réciproque, vérifier que la dérivée ne s'annule pas.

Pour une inégalité, construire une fonction auxiliaire dont on étudie le minimum.

Pour une unicité, chercher une stricte monotonie ou une injectivité.

Pour une existence, chercher un changement de signe ou une limite aux bornes.

Pour une borne, chercher une majoration uniforme de la dérivée.

\end{methodebox}

\subsection*{Grille 5 -- réflexe de rédaction}

\begin{methodebox}

Avant de dériver, écrire le domaine de définition.

Avant d'appliquer Rolle, vérifier la continuité sur le segment.

Avant d'appliquer Rolle, vérifier la dérivabilité sur l'intervalle ouvert.

Avant d'appliquer le TAF, vérifier que les deux bornes appartiennent au même intervalle.

Avant d'utiliser une réciproque, vérifier la bijectivité par continuité et stricte monotonie.

Avant d'utiliser la formule de la réciproque, vérifier que la dérivée ne s'annule pas.

Pour une inégalité, construire une fonction auxiliaire dont on étudie le minimum.

Pour une unicité, chercher une stricte monotonie ou une injectivité.

Pour une existence, chercher un changement de signe ou une limite aux bornes.

Pour une borne, chercher une majoration uniforme de la dérivée.

\end{methodebox}

\subsection*{Grille 6 -- réflexe de rédaction}

\begin{methodebox}

Avant de dériver, écrire le domaine de définition.

Avant d'appliquer Rolle, vérifier la continuité sur le segment.

Avant d'appliquer Rolle, vérifier la dérivabilité sur l'intervalle ouvert.

Avant d'appliquer le TAF, vérifier que les deux bornes appartiennent au même intervalle.

Avant d'utiliser une réciproque, vérifier la bijectivité par continuité et stricte monotonie.

Avant d'utiliser la formule de la réciproque, vérifier que la dérivée ne s'annule pas.

Pour une inégalité, construire une fonction auxiliaire dont on étudie le minimum.

Pour une unicité, chercher une stricte monotonie ou une injectivité.

Pour une existence, chercher un changement de signe ou une limite aux bornes.

Pour une borne, chercher une majoration uniforme de la dérivée.

\end{methodebox}

\subsection*{Grille 7 -- réflexe de rédaction}

\begin{methodebox}

Avant de dériver, écrire le domaine de définition.

Avant d'appliquer Rolle, vérifier la continuité sur le segment.

Avant d'appliquer Rolle, vérifier la dérivabilité sur l'intervalle ouvert.

Avant d'appliquer le TAF, vérifier que les deux bornes appartiennent au même intervalle.

Avant d'utiliser une réciproque, vérifier la bijectivité par continuité et stricte monotonie.

Avant d'utiliser la formule de la réciproque, vérifier que la dérivée ne s'annule pas.

Pour une inégalité, construire une fonction auxiliaire dont on étudie le minimum.

Pour une unicité, chercher une stricte monotonie ou une injectivité.

Pour une existence, chercher un changement de signe ou une limite aux bornes.

Pour une borne, chercher une majoration uniforme de la dérivée.

\end{methodebox}

\subsection*{Grille 8 -- réflexe de rédaction}

\begin{methodebox}

Avant de dériver, écrire le domaine de définition.

Avant d'appliquer Rolle, vérifier la continuité sur le segment.

Avant d'appliquer Rolle, vérifier la dérivabilité sur l'intervalle ouvert.

Avant d'appliquer le TAF, vérifier que les deux bornes appartiennent au même intervalle.

Avant d'utiliser une réciproque, vérifier la bijectivité par continuité et stricte monotonie.

Avant d'utiliser la formule de la réciproque, vérifier que la dérivée ne s'annule pas.

Pour une inégalité, construire une fonction auxiliaire dont on étudie le minimum.

Pour une unicité, chercher une stricte monotonie ou une injectivité.

Pour une existence, chercher un changement de signe ou une limite aux bornes.

Pour une borne, chercher une majoration uniforme de la dérivée.

\end{methodebox}

\subsection*{Grille 9 -- réflexe de rédaction}

\begin{methodebox}

Avant de dériver, écrire le domaine de définition.

Avant d'appliquer Rolle, vérifier la continuité sur le segment.

Avant d'appliquer Rolle, vérifier la dérivabilité sur l'intervalle ouvert.

Avant d'appliquer le TAF, vérifier que les deux bornes appartiennent au même intervalle.

Avant d'utiliser une réciproque, vérifier la bijectivité par continuité et stricte monotonie.

Avant d'utiliser la formule de la réciproque, vérifier que la dérivée ne s'annule pas.

Pour une inégalité, construire une fonction auxiliaire dont on étudie le minimum.

Pour une unicité, chercher une stricte monotonie ou une injectivité.

Pour une existence, chercher un changement de signe ou une limite aux bornes.

Pour une borne, chercher une majoration uniforme de la dérivée.

\end{methodebox}

\subsection*{Grille 10 -- réflexe de rédaction}

\begin{methodebox}

Avant de dériver, écrire le domaine de définition.

Avant d'appliquer Rolle, vérifier la continuité sur le segment.

Avant d'appliquer Rolle, vérifier la dérivabilité sur l'intervalle ouvert.

Avant d'appliquer le TAF, vérifier que les deux bornes appartiennent au même intervalle.

Avant d'utiliser une réciproque, vérifier la bijectivité par continuité et stricte monotonie.

Avant d'utiliser la formule de la réciproque, vérifier que la dérivée ne s'annule pas.

Pour une inégalité, construire une fonction auxiliaire dont on étudie le minimum.

Pour une unicité, chercher une stricte monotonie ou une injectivité.

Pour une existence, chercher un changement de signe ou une limite aux bornes.

Pour une borne, chercher une majoration uniforme de la dérivée.

\end{methodebox}

\subsection*{Grille 11 -- réflexe de rédaction}

\begin{methodebox}

Avant de dériver, écrire le domaine de définition.

Avant d'appliquer Rolle, vérifier la continuité sur le segment.

Avant d'appliquer Rolle, vérifier la dérivabilité sur l'intervalle ouvert.

Avant d'appliquer le TAF, vérifier que les deux bornes appartiennent au même intervalle.

Avant d'utiliser une réciproque, vérifier la bijectivité par continuité et stricte monotonie.

Avant d'utiliser la formule de la réciproque, vérifier que la dérivée ne s'annule pas.

Pour une inégalité, construire une fonction auxiliaire dont on étudie le minimum.

Pour une unicité, chercher une stricte monotonie ou une injectivité.

Pour une existence, chercher un changement de signe ou une limite aux bornes.

Pour une borne, chercher une majoration uniforme de la dérivée.

\end{methodebox}

\subsection*{Grille 12 -- réflexe de rédaction}

\begin{methodebox}

Avant de dériver, écrire le domaine de définition.

Avant d'appliquer Rolle, vérifier la continuité sur le segment.

Avant d'appliquer Rolle, vérifier la dérivabilité sur l'intervalle ouvert.

Avant d'appliquer le TAF, vérifier que les deux bornes appartiennent au même intervalle.

Avant d'utiliser une réciproque, vérifier la bijectivité par continuité et stricte monotonie.

Avant d'utiliser la formule de la réciproque, vérifier que la dérivée ne s'annule pas.

Pour une inégalité, construire une fonction auxiliaire dont on étudie le minimum.

Pour une unicité, chercher une stricte monotonie ou une injectivité.

Pour une existence, chercher un changement de signe ou une limite aux bornes.

Pour une borne, chercher une majoration uniforme de la dérivée.

\end{methodebox}

\subsection*{Grille 13 -- réflexe de rédaction}

\begin{methodebox}

Avant de dériver, écrire le domaine de définition.

Avant d'appliquer Rolle, vérifier la continuité sur le segment.

Avant d'appliquer Rolle, vérifier la dérivabilité sur l'intervalle ouvert.

Avant d'appliquer le TAF, vérifier que les deux bornes appartiennent au même intervalle.

Avant d'utiliser une réciproque, vérifier la bijectivité par continuité et stricte monotonie.

Avant d'utiliser la formule de la réciproque, vérifier que la dérivée ne s'annule pas.

Pour une inégalité, construire une fonction auxiliaire dont on étudie le minimum.

Pour une unicité, chercher une stricte monotonie ou une injectivité.

Pour une existence, chercher un changement de signe ou une limite aux bornes.

Pour une borne, chercher une majoration uniforme de la dérivée.

\end{methodebox}

\subsection*{Grille 14 -- réflexe de rédaction}

\begin{methodebox}

Avant de dériver, écrire le domaine de définition.

Avant d'appliquer Rolle, vérifier la continuité sur le segment.

Avant d'appliquer Rolle, vérifier la dérivabilité sur l'intervalle ouvert.

Avant d'appliquer le TAF, vérifier que les deux bornes appartiennent au même intervalle.

Avant d'utiliser une réciproque, vérifier la bijectivité par continuité et stricte monotonie.

Avant d'utiliser la formule de la réciproque, vérifier que la dérivée ne s'annule pas.

Pour une inégalité, construire une fonction auxiliaire dont on étudie le minimum.

Pour une unicité, chercher une stricte monotonie ou une injectivité.

Pour une existence, chercher un changement de signe ou une limite aux bornes.

Pour une borne, chercher une majoration uniforme de la dérivée.

\end{methodebox}

\subsection*{Grille 15 -- réflexe de rédaction}

\begin{methodebox}

Avant de dériver, écrire le domaine de définition.

Avant d'appliquer Rolle, vérifier la continuité sur le segment.

Avant d'appliquer Rolle, vérifier la dérivabilité sur l'intervalle ouvert.

Avant d'appliquer le TAF, vérifier que les deux bornes appartiennent au même intervalle.

Avant d'utiliser une réciproque, vérifier la bijectivité par continuité et stricte monotonie.

Avant d'utiliser la formule de la réciproque, vérifier que la dérivée ne s'annule pas.

Pour une inégalité, construire une fonction auxiliaire dont on étudie le minimum.

Pour une unicité, chercher une stricte monotonie ou une injectivité.

Pour une existence, chercher un changement de signe ou une limite aux bornes.

Pour une borne, chercher une majoration uniforme de la dérivée.

\end{methodebox}

\subsection*{Grille 16 -- réflexe de rédaction}

\begin{methodebox}

Avant de dériver, écrire le domaine de définition.

Avant d'appliquer Rolle, vérifier la continuité sur le segment.

Avant d'appliquer Rolle, vérifier la dérivabilité sur l'intervalle ouvert.

Avant d'appliquer le TAF, vérifier que les deux bornes appartiennent au même intervalle.

Avant d'utiliser une réciproque, vérifier la bijectivité par continuité et stricte monotonie.

Avant d'utiliser la formule de la réciproque, vérifier que la dérivée ne s'annule pas.

Pour une inégalité, construire une fonction auxiliaire dont on étudie le minimum.

Pour une unicité, chercher une stricte monotonie ou une injectivité.

Pour une existence, chercher un changement de signe ou une limite aux bornes.

Pour une borne, chercher une majoration uniforme de la dérivée.

\end{methodebox}

\subsection*{Grille 17 -- réflexe de rédaction}

\begin{methodebox}

Avant de dériver, écrire le domaine de définition.

Avant d'appliquer Rolle, vérifier la continuité sur le segment.

Avant d'appliquer Rolle, vérifier la dérivabilité sur l'intervalle ouvert.

Avant d'appliquer le TAF, vérifier que les deux bornes appartiennent au même intervalle.

Avant d'utiliser une réciproque, vérifier la bijectivité par continuité et stricte monotonie.

Avant d'utiliser la formule de la réciproque, vérifier que la dérivée ne s'annule pas.

Pour une inégalité, construire une fonction auxiliaire dont on étudie le minimum.

Pour une unicité, chercher une stricte monotonie ou une injectivité.

Pour une existence, chercher un changement de signe ou une limite aux bornes.

Pour une borne, chercher une majoration uniforme de la dérivée.

\end{methodebox}

\subsection*{Grille 18 -- réflexe de rédaction}

\begin{methodebox}

Avant de dériver, écrire le domaine de définition.

Avant d'appliquer Rolle, vérifier la continuité sur le segment.

Avant d'appliquer Rolle, vérifier la dérivabilité sur l'intervalle ouvert.

Avant d'appliquer le TAF, vérifier que les deux bornes appartiennent au même intervalle.

Avant d'utiliser une réciproque, vérifier la bijectivité par continuité et stricte monotonie.

Avant d'utiliser la formule de la réciproque, vérifier que la dérivée ne s'annule pas.

Pour une inégalité, construire une fonction auxiliaire dont on étudie le minimum.

Pour une unicité, chercher une stricte monotonie ou une injectivité.

Pour une existence, chercher un changement de signe ou une limite aux bornes.

Pour une borne, chercher une majoration uniforme de la dérivée.

\end{methodebox}

\subsection*{Grille 19 -- réflexe de rédaction}

\begin{methodebox}

Avant de dériver, écrire le domaine de définition.

Avant d'appliquer Rolle, vérifier la continuité sur le segment.

Avant d'appliquer Rolle, vérifier la dérivabilité sur l'intervalle ouvert.

Avant d'appliquer le TAF, vérifier que les deux bornes appartiennent au même intervalle.

Avant d'utiliser une réciproque, vérifier la bijectivité par continuité et stricte monotonie.

Avant d'utiliser la formule de la réciproque, vérifier que la dérivée ne s'annule pas.

Pour une inégalité, construire une fonction auxiliaire dont on étudie le minimum.

Pour une unicité, chercher une stricte monotonie ou une injectivité.

Pour une existence, chercher un changement de signe ou une limite aux bornes.

Pour une borne, chercher une majoration uniforme de la dérivée.

\end{methodebox}

\subsection*{Grille 20 -- réflexe de rédaction}

\begin{methodebox}

Avant de dériver, écrire le domaine de définition.

Avant d'appliquer Rolle, vérifier la continuité sur le segment.

Avant d'appliquer Rolle, vérifier la dérivabilité sur l'intervalle ouvert.

Avant d'appliquer le TAF, vérifier que les deux bornes appartiennent au même intervalle.

Avant d'utiliser une réciproque, vérifier la bijectivité par continuité et stricte monotonie.

Avant d'utiliser la formule de la réciproque, vérifier que la dérivée ne s'annule pas.

Pour une inégalité, construire une fonction auxiliaire dont on étudie le minimum.

Pour une unicité, chercher une stricte monotonie ou une injectivité.

Pour une existence, chercher un changement de signe ou une limite aux bornes.

Pour une borne, chercher une majoration uniforme de la dérivée.

\end{methodebox}

\subsection*{Grille 21 -- réflexe de rédaction}

\begin{methodebox}

Avant de dériver, écrire le domaine de définition.

Avant d'appliquer Rolle, vérifier la continuité sur le segment.

Avant d'appliquer Rolle, vérifier la dérivabilité sur l'intervalle ouvert.

Avant d'appliquer le TAF, vérifier que les deux bornes appartiennent au même intervalle.

Avant d'utiliser une réciproque, vérifier la bijectivité par continuité et stricte monotonie.

Avant d'utiliser la formule de la réciproque, vérifier que la dérivée ne s'annule pas.

Pour une inégalité, construire une fonction auxiliaire dont on étudie le minimum.

Pour une unicité, chercher une stricte monotonie ou une injectivité.

Pour une existence, chercher un changement de signe ou une limite aux bornes.

Pour une borne, chercher une majoration uniforme de la dérivée.

\end{methodebox}

\subsection*{Grille 22 -- réflexe de rédaction}

\begin{methodebox}

Avant de dériver, écrire le domaine de définition.

Avant d'appliquer Rolle, vérifier la continuité sur le segment.

Avant d'appliquer Rolle, vérifier la dérivabilité sur l'intervalle ouvert.

Avant d'appliquer le TAF, vérifier que les deux bornes appartiennent au même intervalle.

Avant d'utiliser une réciproque, vérifier la bijectivité par continuité et stricte monotonie.

Avant d'utiliser la formule de la réciproque, vérifier que la dérivée ne s'annule pas.

Pour une inégalité, construire une fonction auxiliaire dont on étudie le minimum.

Pour une unicité, chercher une stricte monotonie ou une injectivité.

Pour une existence, chercher un changement de signe ou une limite aux bornes.

Pour une borne, chercher une majoration uniforme de la dérivée.

\end{methodebox}

\subsection*{Grille 23 -- réflexe de rédaction}

\begin{methodebox}

Avant de dériver, écrire le domaine de définition.

Avant d'appliquer Rolle, vérifier la continuité sur le segment.

Avant d'appliquer Rolle, vérifier la dérivabilité sur l'intervalle ouvert.

Avant d'appliquer le TAF, vérifier que les deux bornes appartiennent au même intervalle.

Avant d'utiliser une réciproque, vérifier la bijectivité par continuité et stricte monotonie.

Avant d'utiliser la formule de la réciproque, vérifier que la dérivée ne s'annule pas.

Pour une inégalité, construire une fonction auxiliaire dont on étudie le minimum.

Pour une unicité, chercher une stricte monotonie ou une injectivité.

Pour une existence, chercher un changement de signe ou une limite aux bornes.

Pour une borne, chercher une majoration uniforme de la dérivée.

\end{methodebox}

\subsection*{Grille 24 -- réflexe de rédaction}

\begin{methodebox}

Avant de dériver, écrire le domaine de définition.

Avant d'appliquer Rolle, vérifier la continuité sur le segment.

Avant d'appliquer Rolle, vérifier la dérivabilité sur l'intervalle ouvert.

Avant d'appliquer le TAF, vérifier que les deux bornes appartiennent au même intervalle.

Avant d'utiliser une réciproque, vérifier la bijectivité par continuité et stricte monotonie.

Avant d'utiliser la formule de la réciproque, vérifier que la dérivée ne s'annule pas.

Pour une inégalité, construire une fonction auxiliaire dont on étudie le minimum.

Pour une unicité, chercher une stricte monotonie ou une injectivité.

Pour une existence, chercher un changement de signe ou une limite aux bornes.

Pour une borne, chercher une majoration uniforme de la dérivée.

\end{methodebox}

\subsection*{Grille 25 -- réflexe de rédaction}

\begin{methodebox}

Avant de dériver, écrire le domaine de définition.

Avant d'appliquer Rolle, vérifier la continuité sur le segment.

Avant d'appliquer Rolle, vérifier la dérivabilité sur l'intervalle ouvert.

Avant d'appliquer le TAF, vérifier que les deux bornes appartiennent au même intervalle.

Avant d'utiliser une réciproque, vérifier la bijectivité par continuité et stricte monotonie.

Avant d'utiliser la formule de la réciproque, vérifier que la dérivée ne s'annule pas.

Pour une inégalité, construire une fonction auxiliaire dont on étudie le minimum.

Pour une unicité, chercher une stricte monotonie ou une injectivité.

Pour une existence, chercher un changement de signe ou une limite aux bornes.

Pour une borne, chercher une majoration uniforme de la dérivée.

\end{methodebox}

\subsection*{Grille 26 -- réflexe de rédaction}

\begin{methodebox}

Avant de dériver, écrire le domaine de définition.

Avant d'appliquer Rolle, vérifier la continuité sur le segment.

Avant d'appliquer Rolle, vérifier la dérivabilité sur l'intervalle ouvert.

Avant d'appliquer le TAF, vérifier que les deux bornes appartiennent au même intervalle.

Avant d'utiliser une réciproque, vérifier la bijectivité par continuité et stricte monotonie.

Avant d'utiliser la formule de la réciproque, vérifier que la dérivée ne s'annule pas.

Pour une inégalité, construire une fonction auxiliaire dont on étudie le minimum.

Pour une unicité, chercher une stricte monotonie ou une injectivité.

Pour une existence, chercher un changement de signe ou une limite aux bornes.

Pour une borne, chercher une majoration uniforme de la dérivée.

\end{methodebox}

\subsection*{Grille 27 -- réflexe de rédaction}

\begin{methodebox}

Avant de dériver, écrire le domaine de définition.

Avant d'appliquer Rolle, vérifier la continuité sur le segment.

Avant d'appliquer Rolle, vérifier la dérivabilité sur l'intervalle ouvert.

Avant d'appliquer le TAF, vérifier que les deux bornes appartiennent au même intervalle.

Avant d'utiliser une réciproque, vérifier la bijectivité par continuité et stricte monotonie.

Avant d'utiliser la formule de la réciproque, vérifier que la dérivée ne s'annule pas.

Pour une inégalité, construire une fonction auxiliaire dont on étudie le minimum.

Pour une unicité, chercher une stricte monotonie ou une injectivité.

Pour une existence, chercher un changement de signe ou une limite aux bornes.

Pour une borne, chercher une majoration uniforme de la dérivée.

\end{methodebox}

\subsection*{Grille 28 -- réflexe de rédaction}

\begin{methodebox}

Avant de dériver, écrire le domaine de définition.

Avant d'appliquer Rolle, vérifier la continuité sur le segment.

Avant d'appliquer Rolle, vérifier la dérivabilité sur l'intervalle ouvert.

Avant d'appliquer le TAF, vérifier que les deux bornes appartiennent au même intervalle.

Avant d'utiliser une réciproque, vérifier la bijectivité par continuité et stricte monotonie.

Avant d'utiliser la formule de la réciproque, vérifier que la dérivée ne s'annule pas.

Pour une inégalité, construire une fonction auxiliaire dont on étudie le minimum.

Pour une unicité, chercher une stricte monotonie ou une injectivité.

Pour une existence, chercher un changement de signe ou une limite aux bornes.

Pour une borne, chercher une majoration uniforme de la dérivée.

\end{methodebox}

\subsection*{Grille 29 -- réflexe de rédaction}

\begin{methodebox}

Avant de dériver, écrire le domaine de définition.

Avant d'appliquer Rolle, vérifier la continuité sur le segment.

Avant d'appliquer Rolle, vérifier la dérivabilité sur l'intervalle ouvert.

Avant d'appliquer le TAF, vérifier que les deux bornes appartiennent au même intervalle.

Avant d'utiliser une réciproque, vérifier la bijectivité par continuité et stricte monotonie.

Avant d'utiliser la formule de la réciproque, vérifier que la dérivée ne s'annule pas.

Pour une inégalité, construire une fonction auxiliaire dont on étudie le minimum.

Pour une unicité, chercher une stricte monotonie ou une injectivité.

Pour une existence, chercher un changement de signe ou une limite aux bornes.

Pour une borne, chercher une majoration uniforme de la dérivée.

\end{methodebox}

\subsection*{Grille 30 -- réflexe de rédaction}

\begin{methodebox}

Avant de dériver, écrire le domaine de définition.

Avant d'appliquer Rolle, vérifier la continuité sur le segment.

Avant d'appliquer Rolle, vérifier la dérivabilité sur l'intervalle ouvert.

Avant d'appliquer le TAF, vérifier que les deux bornes appartiennent au même intervalle.

Avant d'utiliser une réciproque, vérifier la bijectivité par continuité et stricte monotonie.

Avant d'utiliser la formule de la réciproque, vérifier que la dérivée ne s'annule pas.

Pour une inégalité, construire une fonction auxiliaire dont on étudie le minimum.

Pour une unicité, chercher une stricte monotonie ou une injectivité.

Pour une existence, chercher un changement de signe ou une limite aux bornes.

Pour une borne, chercher une majoration uniforme de la dérivée.

\end{methodebox}

\subsection*{Grille 31 -- réflexe de rédaction}

\begin{methodebox}

Avant de dériver, écrire le domaine de définition.

Avant d'appliquer Rolle, vérifier la continuité sur le segment.

Avant d'appliquer Rolle, vérifier la dérivabilité sur l'intervalle ouvert.

Avant d'appliquer le TAF, vérifier que les deux bornes appartiennent au même intervalle.

Avant d'utiliser une réciproque, vérifier la bijectivité par continuité et stricte monotonie.

Avant d'utiliser la formule de la réciproque, vérifier que la dérivée ne s'annule pas.

Pour une inégalité, construire une fonction auxiliaire dont on étudie le minimum.

Pour une unicité, chercher une stricte monotonie ou une injectivité.

Pour une existence, chercher un changement de signe ou une limite aux bornes.

Pour une borne, chercher une majoration uniforme de la dérivée.

\end{methodebox}

\subsection*{Grille 32 -- réflexe de rédaction}

\begin{methodebox}

Avant de dériver, écrire le domaine de définition.

Avant d'appliquer Rolle, vérifier la continuité sur le segment.

Avant d'appliquer Rolle, vérifier la dérivabilité sur l'intervalle ouvert.

Avant d'appliquer le TAF, vérifier que les deux bornes appartiennent au même intervalle.

Avant d'utiliser une réciproque, vérifier la bijectivité par continuité et stricte monotonie.

Avant d'utiliser la formule de la réciproque, vérifier que la dérivée ne s'annule pas.

Pour une inégalité, construire une fonction auxiliaire dont on étudie le minimum.

Pour une unicité, chercher une stricte monotonie ou une injectivité.

Pour une existence, chercher un changement de signe ou une limite aux bornes.

Pour une borne, chercher une majoration uniforme de la dérivée.

\end{methodebox}

\subsection*{Grille 33 -- réflexe de rédaction}

\begin{methodebox}

Avant de dériver, écrire le domaine de définition.

Avant d'appliquer Rolle, vérifier la continuité sur le segment.

Avant d'appliquer Rolle, vérifier la dérivabilité sur l'intervalle ouvert.

Avant d'appliquer le TAF, vérifier que les deux bornes appartiennent au même intervalle.

Avant d'utiliser une réciproque, vérifier la bijectivité par continuité et stricte monotonie.

Avant d'utiliser la formule de la réciproque, vérifier que la dérivée ne s'annule pas.

Pour une inégalité, construire une fonction auxiliaire dont on étudie le minimum.

Pour une unicité, chercher une stricte monotonie ou une injectivité.

Pour une existence, chercher un changement de signe ou une limite aux bornes.

Pour une borne, chercher une majoration uniforme de la dérivée.

\end{methodebox}

\subsection*{Grille 34 -- réflexe de rédaction}

\begin{methodebox}

Avant de dériver, écrire le domaine de définition.

Avant d'appliquer Rolle, vérifier la continuité sur le segment.

Avant d'appliquer Rolle, vérifier la dérivabilité sur l'intervalle ouvert.

Avant d'appliquer le TAF, vérifier que les deux bornes appartiennent au même intervalle.

Avant d'utiliser une réciproque, vérifier la bijectivité par continuité et stricte monotonie.

Avant d'utiliser la formule de la réciproque, vérifier que la dérivée ne s'annule pas.

Pour une inégalité, construire une fonction auxiliaire dont on étudie le minimum.

Pour une unicité, chercher une stricte monotonie ou une injectivité.

Pour une existence, chercher un changement de signe ou une limite aux bornes.

Pour une borne, chercher une majoration uniforme de la dérivée.

\end{methodebox}

\subsection*{Grille 35 -- réflexe de rédaction}

\begin{methodebox}

Avant de dériver, écrire le domaine de définition.

Avant d'appliquer Rolle, vérifier la continuité sur le segment.

Avant d'appliquer Rolle, vérifier la dérivabilité sur l'intervalle ouvert.

Avant d'appliquer le TAF, vérifier que les deux bornes appartiennent au même intervalle.

Avant d'utiliser une réciproque, vérifier la bijectivité par continuité et stricte monotonie.

Avant d'utiliser la formule de la réciproque, vérifier que la dérivée ne s'annule pas.

Pour une inégalité, construire une fonction auxiliaire dont on étudie le minimum.

Pour une unicité, chercher une stricte monotonie ou une injectivité.

Pour une existence, chercher un changement de signe ou une limite aux bornes.

Pour une borne, chercher une majoration uniforme de la dérivée.

\end{methodebox}

\subsection*{Grille 36 -- réflexe de rédaction}

\begin{methodebox}

Avant de dériver, écrire le domaine de définition.

Avant d'appliquer Rolle, vérifier la continuité sur le segment.

Avant d'appliquer Rolle, vérifier la dérivabilité sur l'intervalle ouvert.

Avant d'appliquer le TAF, vérifier que les deux bornes appartiennent au même intervalle.

Avant d'utiliser une réciproque, vérifier la bijectivité par continuité et stricte monotonie.

Avant d'utiliser la formule de la réciproque, vérifier que la dérivée ne s'annule pas.

Pour une inégalité, construire une fonction auxiliaire dont on étudie le minimum.

Pour une unicité, chercher une stricte monotonie ou une injectivité.

Pour une existence, chercher un changement de signe ou une limite aux bornes.

Pour une borne, chercher une majoration uniforme de la dérivée.

\end{methodebox}

\subsection*{Grille 37 -- réflexe de rédaction}

\begin{methodebox}

Avant de dériver, écrire le domaine de définition.

Avant d'appliquer Rolle, vérifier la continuité sur le segment.

Avant d'appliquer Rolle, vérifier la dérivabilité sur l'intervalle ouvert.

Avant d'appliquer le TAF, vérifier que les deux bornes appartiennent au même intervalle.

Avant d'utiliser une réciproque, vérifier la bijectivité par continuité et stricte monotonie.

Avant d'utiliser la formule de la réciproque, vérifier que la dérivée ne s'annule pas.

Pour une inégalité, construire une fonction auxiliaire dont on étudie le minimum.

Pour une unicité, chercher une stricte monotonie ou une injectivité.

Pour une existence, chercher un changement de signe ou une limite aux bornes.

Pour une borne, chercher une majoration uniforme de la dérivée.

\end{methodebox}

\subsection*{Grille 38 -- réflexe de rédaction}

\begin{methodebox}

Avant de dériver, écrire le domaine de définition.

Avant d'appliquer Rolle, vérifier la continuité sur le segment.

Avant d'appliquer Rolle, vérifier la dérivabilité sur l'intervalle ouvert.

Avant d'appliquer le TAF, vérifier que les deux bornes appartiennent au même intervalle.

Avant d'utiliser une réciproque, vérifier la bijectivité par continuité et stricte monotonie.

Avant d'utiliser la formule de la réciproque, vérifier que la dérivée ne s'annule pas.

Pour une inégalité, construire une fonction auxiliaire dont on étudie le minimum.

Pour une unicité, chercher une stricte monotonie ou une injectivité.

Pour une existence, chercher un changement de signe ou une limite aux bornes.

Pour une borne, chercher une majoration uniforme de la dérivée.

\end{methodebox}

\subsection*{Grille 39 -- réflexe de rédaction}

\begin{methodebox}

Avant de dériver, écrire le domaine de définition.

Avant d'appliquer Rolle, vérifier la continuité sur le segment.

Avant d'appliquer Rolle, vérifier la dérivabilité sur l'intervalle ouvert.

Avant d'appliquer le TAF, vérifier que les deux bornes appartiennent au même intervalle.

Avant d'utiliser une réciproque, vérifier la bijectivité par continuité et stricte monotonie.

Avant d'utiliser la formule de la réciproque, vérifier que la dérivée ne s'annule pas.

Pour une inégalité, construire une fonction auxiliaire dont on étudie le minimum.

Pour une unicité, chercher une stricte monotonie ou une injectivité.

Pour une existence, chercher un changement de signe ou une limite aux bornes.

Pour une borne, chercher une majoration uniforme de la dérivée.

\end{methodebox}

\subsection*{Grille 40 -- réflexe de rédaction}

\begin{methodebox}

Avant de dériver, écrire le domaine de définition.

Avant d'appliquer Rolle, vérifier la continuité sur le segment.

Avant d'appliquer Rolle, vérifier la dérivabilité sur l'intervalle ouvert.

Avant d'appliquer le TAF, vérifier que les deux bornes appartiennent au même intervalle.

Avant d'utiliser une réciproque, vérifier la bijectivité par continuité et stricte monotonie.

Avant d'utiliser la formule de la réciproque, vérifier que la dérivée ne s'annule pas.

Pour une inégalité, construire une fonction auxiliaire dont on étudie le minimum.

Pour une unicité, chercher une stricte monotonie ou une injectivité.

Pour une existence, chercher un changement de signe ou une limite aux bornes.

Pour une borne, chercher une majoration uniforme de la dérivée.

\end{methodebox}

\subsection*{Grille 41 -- réflexe de rédaction}

\begin{methodebox}

Avant de dériver, écrire le domaine de définition.

Avant d'appliquer Rolle, vérifier la continuité sur le segment.

Avant d'appliquer Rolle, vérifier la dérivabilité sur l'intervalle ouvert.

Avant d'appliquer le TAF, vérifier que les deux bornes appartiennent au même intervalle.

Avant d'utiliser une réciproque, vérifier la bijectivité par continuité et stricte monotonie.

Avant d'utiliser la formule de la réciproque, vérifier que la dérivée ne s'annule pas.

Pour une inégalité, construire une fonction auxiliaire dont on étudie le minimum.

Pour une unicité, chercher une stricte monotonie ou une injectivité.

Pour une existence, chercher un changement de signe ou une limite aux bornes.

Pour une borne, chercher une majoration uniforme de la dérivée.

\end{methodebox}

\subsection*{Grille 42 -- réflexe de rédaction}

\begin{methodebox}

Avant de dériver, écrire le domaine de définition.

Avant d'appliquer Rolle, vérifier la continuité sur le segment.

Avant d'appliquer Rolle, vérifier la dérivabilité sur l'intervalle ouvert.

Avant d'appliquer le TAF, vérifier que les deux bornes appartiennent au même intervalle.

Avant d'utiliser une réciproque, vérifier la bijectivité par continuité et stricte monotonie.

Avant d'utiliser la formule de la réciproque, vérifier que la dérivée ne s'annule pas.

Pour une inégalité, construire une fonction auxiliaire dont on étudie le minimum.

Pour une unicité, chercher une stricte monotonie ou une injectivité.

Pour une existence, chercher un changement de signe ou une limite aux bornes.

Pour une borne, chercher une majoration uniforme de la dérivée.

\end{methodebox}

\subsection*{Grille 43 -- réflexe de rédaction}

\begin{methodebox}

Avant de dériver, écrire le domaine de définition.

Avant d'appliquer Rolle, vérifier la continuité sur le segment.

Avant d'appliquer Rolle, vérifier la dérivabilité sur l'intervalle ouvert.

Avant d'appliquer le TAF, vérifier que les deux bornes appartiennent au même intervalle.

Avant d'utiliser une réciproque, vérifier la bijectivité par continuité et stricte monotonie.

Avant d'utiliser la formule de la réciproque, vérifier que la dérivée ne s'annule pas.

Pour une inégalité, construire une fonction auxiliaire dont on étudie le minimum.

Pour une unicité, chercher une stricte monotonie ou une injectivité.

Pour une existence, chercher un changement de signe ou une limite aux bornes.

Pour une borne, chercher une majoration uniforme de la dérivée.

\end{methodebox}

\subsection*{Grille 44 -- réflexe de rédaction}

\begin{methodebox}

Avant de dériver, écrire le domaine de définition.

Avant d'appliquer Rolle, vérifier la continuité sur le segment.

Avant d'appliquer Rolle, vérifier la dérivabilité sur l'intervalle ouvert.

Avant d'appliquer le TAF, vérifier que les deux bornes appartiennent au même intervalle.

Avant d'utiliser une réciproque, vérifier la bijectivité par continuité et stricte monotonie.

Avant d'utiliser la formule de la réciproque, vérifier que la dérivée ne s'annule pas.

Pour une inégalité, construire une fonction auxiliaire dont on étudie le minimum.

Pour une unicité, chercher une stricte monotonie ou une injectivité.

Pour une existence, chercher un changement de signe ou une limite aux bornes.

Pour une borne, chercher une majoration uniforme de la dérivée.

\end{methodebox}

\subsection*{Grille 45 -- réflexe de rédaction}

\begin{methodebox}

Avant de dériver, écrire le domaine de définition.

Avant d'appliquer Rolle, vérifier la continuité sur le segment.

Avant d'appliquer Rolle, vérifier la dérivabilité sur l'intervalle ouvert.

Avant d'appliquer le TAF, vérifier que les deux bornes appartiennent au même intervalle.

Avant d'utiliser une réciproque, vérifier la bijectivité par continuité et stricte monotonie.

Avant d'utiliser la formule de la réciproque, vérifier que la dérivée ne s'annule pas.

Pour une inégalité, construire une fonction auxiliaire dont on étudie le minimum.

Pour une unicité, chercher une stricte monotonie ou une injectivité.

Pour une existence, chercher un changement de signe ou une limite aux bornes.

Pour une borne, chercher une majoration uniforme de la dérivée.

\end{methodebox}

\subsection*{Grille 46 -- réflexe de rédaction}

\begin{methodebox}

Avant de dériver, écrire le domaine de définition.

Avant d'appliquer Rolle, vérifier la continuité sur le segment.

Avant d'appliquer Rolle, vérifier la dérivabilité sur l'intervalle ouvert.

Avant d'appliquer le TAF, vérifier que les deux bornes appartiennent au même intervalle.

Avant d'utiliser une réciproque, vérifier la bijectivité par continuité et stricte monotonie.

Avant d'utiliser la formule de la réciproque, vérifier que la dérivée ne s'annule pas.

Pour une inégalité, construire une fonction auxiliaire dont on étudie le minimum.

Pour une unicité, chercher une stricte monotonie ou une injectivité.

Pour une existence, chercher un changement de signe ou une limite aux bornes.

Pour une borne, chercher une majoration uniforme de la dérivée.

\end{methodebox}

\subsection*{Grille 47 -- réflexe de rédaction}

\begin{methodebox}

Avant de dériver, écrire le domaine de définition.

Avant d'appliquer Rolle, vérifier la continuité sur le segment.

Avant d'appliquer Rolle, vérifier la dérivabilité sur l'intervalle ouvert.

Avant d'appliquer le TAF, vérifier que les deux bornes appartiennent au même intervalle.

Avant d'utiliser une réciproque, vérifier la bijectivité par continuité et stricte monotonie.

Avant d'utiliser la formule de la réciproque, vérifier que la dérivée ne s'annule pas.

Pour une inégalité, construire une fonction auxiliaire dont on étudie le minimum.

Pour une unicité, chercher une stricte monotonie ou une injectivité.

Pour une existence, chercher un changement de signe ou une limite aux bornes.

Pour une borne, chercher une majoration uniforme de la dérivée.

\end{methodebox}

\subsection*{Grille 48 -- réflexe de rédaction}

\begin{methodebox}

Avant de dériver, écrire le domaine de définition.

Avant d'appliquer Rolle, vérifier la continuité sur le segment.

Avant d'appliquer Rolle, vérifier la dérivabilité sur l'intervalle ouvert.

Avant d'appliquer le TAF, vérifier que les deux bornes appartiennent au même intervalle.

Avant d'utiliser une réciproque, vérifier la bijectivité par continuité et stricte monotonie.

Avant d'utiliser la formule de la réciproque, vérifier que la dérivée ne s'annule pas.

Pour une inégalité, construire une fonction auxiliaire dont on étudie le minimum.

Pour une unicité, chercher une stricte monotonie ou une injectivité.

Pour une existence, chercher un changement de signe ou une limite aux bornes.

Pour une borne, chercher une majoration uniforme de la dérivée.

\end{methodebox}

\subsection*{Grille 49 -- réflexe de rédaction}

\begin{methodebox}

Avant de dériver, écrire le domaine de définition.

Avant d'appliquer Rolle, vérifier la continuité sur le segment.

Avant d'appliquer Rolle, vérifier la dérivabilité sur l'intervalle ouvert.

Avant d'appliquer le TAF, vérifier que les deux bornes appartiennent au même intervalle.

Avant d'utiliser une réciproque, vérifier la bijectivité par continuité et stricte monotonie.

Avant d'utiliser la formule de la réciproque, vérifier que la dérivée ne s'annule pas.

Pour une inégalité, construire une fonction auxiliaire dont on étudie le minimum.

Pour une unicité, chercher une stricte monotonie ou une injectivité.

Pour une existence, chercher un changement de signe ou une limite aux bornes.

Pour une borne, chercher une majoration uniforme de la dérivée.

\end{methodebox}

\subsection*{Grille 50 -- réflexe de rédaction}

\begin{methodebox}

Avant de dériver, écrire le domaine de définition.

Avant d'appliquer Rolle, vérifier la continuité sur le segment.

Avant d'appliquer Rolle, vérifier la dérivabilité sur l'intervalle ouvert.

Avant d'appliquer le TAF, vérifier que les deux bornes appartiennent au même intervalle.

Avant d'utiliser une réciproque, vérifier la bijectivité par continuité et stricte monotonie.

Avant d'utiliser la formule de la réciproque, vérifier que la dérivée ne s'annule pas.

Pour une inégalité, construire une fonction auxiliaire dont on étudie le minimum.

Pour une unicité, chercher une stricte monotonie ou une injectivité.

Pour une existence, chercher un changement de signe ou une limite aux bornes.

Pour une borne, chercher une majoration uniforme de la dérivée.

\end{methodebox}

\subsection*{Grille 51 -- réflexe de rédaction}

\begin{methodebox}

Avant de dériver, écrire le domaine de définition.

Avant d'appliquer Rolle, vérifier la continuité sur le segment.

Avant d'appliquer Rolle, vérifier la dérivabilité sur l'intervalle ouvert.

Avant d'appliquer le TAF, vérifier que les deux bornes appartiennent au même intervalle.

Avant d'utiliser une réciproque, vérifier la bijectivité par continuité et stricte monotonie.

Avant d'utiliser la formule de la réciproque, vérifier que la dérivée ne s'annule pas.

Pour une inégalité, construire une fonction auxiliaire dont on étudie le minimum.

Pour une unicité, chercher une stricte monotonie ou une injectivité.

Pour une existence, chercher un changement de signe ou une limite aux bornes.

Pour une borne, chercher une majoration uniforme de la dérivée.

\end{methodebox}

\subsection*{Grille 52 -- réflexe de rédaction}

\begin{methodebox}

Avant de dériver, écrire le domaine de définition.

Avant d'appliquer Rolle, vérifier la continuité sur le segment.

Avant d'appliquer Rolle, vérifier la dérivabilité sur l'intervalle ouvert.

Avant d'appliquer le TAF, vérifier que les deux bornes appartiennent au même intervalle.

Avant d'utiliser une réciproque, vérifier la bijectivité par continuité et stricte monotonie.

Avant d'utiliser la formule de la réciproque, vérifier que la dérivée ne s'annule pas.

Pour une inégalité, construire une fonction auxiliaire dont on étudie le minimum.

Pour une unicité, chercher une stricte monotonie ou une injectivité.

Pour une existence, chercher un changement de signe ou une limite aux bornes.

Pour une borne, chercher une majoration uniforme de la dérivée.

\end{methodebox}

\subsection*{Grille 53 -- réflexe de rédaction}

\begin{methodebox}

Avant de dériver, écrire le domaine de définition.

Avant d'appliquer Rolle, vérifier la continuité sur le segment.

Avant d'appliquer Rolle, vérifier la dérivabilité sur l'intervalle ouvert.

Avant d'appliquer le TAF, vérifier que les deux bornes appartiennent au même intervalle.

Avant d'utiliser une réciproque, vérifier la bijectivité par continuité et stricte monotonie.

Avant d'utiliser la formule de la réciproque, vérifier que la dérivée ne s'annule pas.

Pour une inégalité, construire une fonction auxiliaire dont on étudie le minimum.

Pour une unicité, chercher une stricte monotonie ou une injectivité.

Pour une existence, chercher un changement de signe ou une limite aux bornes.

Pour une borne, chercher une majoration uniforme de la dérivée.

\end{methodebox}

\subsection*{Grille 54 -- réflexe de rédaction}

\begin{methodebox}

Avant de dériver, écrire le domaine de définition.

Avant d'appliquer Rolle, vérifier la continuité sur le segment.

Avant d'appliquer Rolle, vérifier la dérivabilité sur l'intervalle ouvert.

Avant d'appliquer le TAF, vérifier que les deux bornes appartiennent au même intervalle.

Avant d'utiliser une réciproque, vérifier la bijectivité par continuité et stricte monotonie.

Avant d'utiliser la formule de la réciproque, vérifier que la dérivée ne s'annule pas.

Pour une inégalité, construire une fonction auxiliaire dont on étudie le minimum.

Pour une unicité, chercher une stricte monotonie ou une injectivité.

Pour une existence, chercher un changement de signe ou une limite aux bornes.

Pour une borne, chercher une majoration uniforme de la dérivée.

\end{methodebox}

\subsection*{Grille 55 -- réflexe de rédaction}

\begin{methodebox}

Avant de dériver, écrire le domaine de définition.

Avant d'appliquer Rolle, vérifier la continuité sur le segment.

Avant d'appliquer Rolle, vérifier la dérivabilité sur l'intervalle ouvert.

Avant d'appliquer le TAF, vérifier que les deux bornes appartiennent au même intervalle.

Avant d'utiliser une réciproque, vérifier la bijectivité par continuité et stricte monotonie.

Avant d'utiliser la formule de la réciproque, vérifier que la dérivée ne s'annule pas.

Pour une inégalité, construire une fonction auxiliaire dont on étudie le minimum.

Pour une unicité, chercher une stricte monotonie ou une injectivité.

Pour une existence, chercher un changement de signe ou une limite aux bornes.

Pour une borne, chercher une majoration uniforme de la dérivée.

\end{methodebox}

\subsection*{Grille 56 -- réflexe de rédaction}

\begin{methodebox}

Avant de dériver, écrire le domaine de définition.

Avant d'appliquer Rolle, vérifier la continuité sur le segment.

Avant d'appliquer Rolle, vérifier la dérivabilité sur l'intervalle ouvert.

Avant d'appliquer le TAF, vérifier que les deux bornes appartiennent au même intervalle.

Avant d'utiliser une réciproque, vérifier la bijectivité par continuité et stricte monotonie.

Avant d'utiliser la formule de la réciproque, vérifier que la dérivée ne s'annule pas.

Pour une inégalité, construire une fonction auxiliaire dont on étudie le minimum.

Pour une unicité, chercher une stricte monotonie ou une injectivité.

Pour une existence, chercher un changement de signe ou une limite aux bornes.

Pour une borne, chercher une majoration uniforme de la dérivée.

\end{methodebox}

\subsection*{Grille 57 -- réflexe de rédaction}

\begin{methodebox}

Avant de dériver, écrire le domaine de définition.

Avant d'appliquer Rolle, vérifier la continuité sur le segment.

Avant d'appliquer Rolle, vérifier la dérivabilité sur l'intervalle ouvert.

Avant d'appliquer le TAF, vérifier que les deux bornes appartiennent au même intervalle.

Avant d'utiliser une réciproque, vérifier la bijectivité par continuité et stricte monotonie.

Avant d'utiliser la formule de la réciproque, vérifier que la dérivée ne s'annule pas.

Pour une inégalité, construire une fonction auxiliaire dont on étudie le minimum.

Pour une unicité, chercher une stricte monotonie ou une injectivité.

Pour une existence, chercher un changement de signe ou une limite aux bornes.

Pour une borne, chercher une majoration uniforme de la dérivée.

\end{methodebox}

\subsection*{Grille 58 -- réflexe de rédaction}

\begin{methodebox}

Avant de dériver, écrire le domaine de définition.

Avant d'appliquer Rolle, vérifier la continuité sur le segment.

Avant d'appliquer Rolle, vérifier la dérivabilité sur l'intervalle ouvert.

Avant d'appliquer le TAF, vérifier que les deux bornes appartiennent au même intervalle.

Avant d'utiliser une réciproque, vérifier la bijectivité par continuité et stricte monotonie.

Avant d'utiliser la formule de la réciproque, vérifier que la dérivée ne s'annule pas.

Pour une inégalité, construire une fonction auxiliaire dont on étudie le minimum.

Pour une unicité, chercher une stricte monotonie ou une injectivité.

Pour une existence, chercher un changement de signe ou une limite aux bornes.

Pour une borne, chercher une majoration uniforme de la dérivée.

\end{methodebox}

\subsection*{Grille 59 -- réflexe de rédaction}

\begin{methodebox}

Avant de dériver, écrire le domaine de définition.

Avant d'appliquer Rolle, vérifier la continuité sur le segment.

Avant d'appliquer Rolle, vérifier la dérivabilité sur l'intervalle ouvert.

Avant d'appliquer le TAF, vérifier que les deux bornes appartiennent au même intervalle.

Avant d'utiliser une réciproque, vérifier la bijectivité par continuité et stricte monotonie.

Avant d'utiliser la formule de la réciproque, vérifier que la dérivée ne s'annule pas.

Pour une inégalité, construire une fonction auxiliaire dont on étudie le minimum.

Pour une unicité, chercher une stricte monotonie ou une injectivité.

Pour une existence, chercher un changement de signe ou une limite aux bornes.

Pour une borne, chercher une majoration uniforme de la dérivée.

\end{methodebox}

\subsection*{Grille 60 -- réflexe de rédaction}

\begin{methodebox}

Avant de dériver, écrire le domaine de définition.

Avant d'appliquer Rolle, vérifier la continuité sur le segment.

Avant d'appliquer Rolle, vérifier la dérivabilité sur l'intervalle ouvert.

Avant d'appliquer le TAF, vérifier que les deux bornes appartiennent au même intervalle.

Avant d'utiliser une réciproque, vérifier la bijectivité par continuité et stricte monotonie.

Avant d'utiliser la formule de la réciproque, vérifier que la dérivée ne s'annule pas.

Pour une inégalité, construire une fonction auxiliaire dont on étudie le minimum.

Pour une unicité, chercher une stricte monotonie ou une injectivité.

Pour une existence, chercher un changement de signe ou une limite aux bornes.

Pour une borne, chercher une majoration uniforme de la dérivée.

\end{methodebox}

\subsection*{Grille 61 -- réflexe de rédaction}

\begin{methodebox}

Avant de dériver, écrire le domaine de définition.

Avant d'appliquer Rolle, vérifier la continuité sur le segment.

Avant d'appliquer Rolle, vérifier la dérivabilité sur l'intervalle ouvert.

Avant d'appliquer le TAF, vérifier que les deux bornes appartiennent au même intervalle.

Avant d'utiliser une réciproque, vérifier la bijectivité par continuité et stricte monotonie.

Avant d'utiliser la formule de la réciproque, vérifier que la dérivée ne s'annule pas.

Pour une inégalité, construire une fonction auxiliaire dont on étudie le minimum.

Pour une unicité, chercher une stricte monotonie ou une injectivité.

Pour une existence, chercher un changement de signe ou une limite aux bornes.

Pour une borne, chercher une majoration uniforme de la dérivée.

\end{methodebox}

\subsection*{Grille 62 -- réflexe de rédaction}

\begin{methodebox}

Avant de dériver, écrire le domaine de définition.

Avant d'appliquer Rolle, vérifier la continuité sur le segment.

Avant d'appliquer Rolle, vérifier la dérivabilité sur l'intervalle ouvert.

Avant d'appliquer le TAF, vérifier que les deux bornes appartiennent au même intervalle.

Avant d'utiliser une réciproque, vérifier la bijectivité par continuité et stricte monotonie.

Avant d'utiliser la formule de la réciproque, vérifier que la dérivée ne s'annule pas.

Pour une inégalité, construire une fonction auxiliaire dont on étudie le minimum.

Pour une unicité, chercher une stricte monotonie ou une injectivité.

Pour une existence, chercher un changement de signe ou une limite aux bornes.

Pour une borne, chercher une majoration uniforme de la dérivée.

\end{methodebox}

\subsection*{Grille 63 -- réflexe de rédaction}

\begin{methodebox}

Avant de dériver, écrire le domaine de définition.

Avant d'appliquer Rolle, vérifier la continuité sur le segment.

Avant d'appliquer Rolle, vérifier la dérivabilité sur l'intervalle ouvert.

Avant d'appliquer le TAF, vérifier que les deux bornes appartiennent au même intervalle.

Avant d'utiliser une réciproque, vérifier la bijectivité par continuité et stricte monotonie.

Avant d'utiliser la formule de la réciproque, vérifier que la dérivée ne s'annule pas.

Pour une inégalité, construire une fonction auxiliaire dont on étudie le minimum.

Pour une unicité, chercher une stricte monotonie ou une injectivité.

Pour une existence, chercher un changement de signe ou une limite aux bornes.

Pour une borne, chercher une majoration uniforme de la dérivée.

\end{methodebox}

\subsection*{Grille 64 -- réflexe de rédaction}

\begin{methodebox}

Avant de dériver, écrire le domaine de définition.

Avant d'appliquer Rolle, vérifier la continuité sur le segment.

Avant d'appliquer Rolle, vérifier la dérivabilité sur l'intervalle ouvert.

Avant d'appliquer le TAF, vérifier que les deux bornes appartiennent au même intervalle.

Avant d'utiliser une réciproque, vérifier la bijectivité par continuité et stricte monotonie.

Avant d'utiliser la formule de la réciproque, vérifier que la dérivée ne s'annule pas.

Pour une inégalité, construire une fonction auxiliaire dont on étudie le minimum.

Pour une unicité, chercher une stricte monotonie ou une injectivité.

Pour une existence, chercher un changement de signe ou une limite aux bornes.

Pour une borne, chercher une majoration uniforme de la dérivée.

\end{methodebox}

\subsection*{Grille 65 -- réflexe de rédaction}

\begin{methodebox}

Avant de dériver, écrire le domaine de définition.

Avant d'appliquer Rolle, vérifier la continuité sur le segment.

Avant d'appliquer Rolle, vérifier la dérivabilité sur l'intervalle ouvert.

Avant d'appliquer le TAF, vérifier que les deux bornes appartiennent au même intervalle.

Avant d'utiliser une réciproque, vérifier la bijectivité par continuité et stricte monotonie.

Avant d'utiliser la formule de la réciproque, vérifier que la dérivée ne s'annule pas.

Pour une inégalité, construire une fonction auxiliaire dont on étudie le minimum.

Pour une unicité, chercher une stricte monotonie ou une injectivité.

Pour une existence, chercher un changement de signe ou une limite aux bornes.

Pour une borne, chercher une majoration uniforme de la dérivée.

\end{methodebox}

\subsection*{Grille 66 -- réflexe de rédaction}

\begin{methodebox}

Avant de dériver, écrire le domaine de définition.

Avant d'appliquer Rolle, vérifier la continuité sur le segment.

Avant d'appliquer Rolle, vérifier la dérivabilité sur l'intervalle ouvert.

Avant d'appliquer le TAF, vérifier que les deux bornes appartiennent au même intervalle.

Avant d'utiliser une réciproque, vérifier la bijectivité par continuité et stricte monotonie.

Avant d'utiliser la formule de la réciproque, vérifier que la dérivée ne s'annule pas.

Pour une inégalité, construire une fonction auxiliaire dont on étudie le minimum.

Pour une unicité, chercher une stricte monotonie ou une injectivité.

Pour une existence, chercher un changement de signe ou une limite aux bornes.

Pour une borne, chercher une majoration uniforme de la dérivée.

\end{methodebox}

\subsection*{Grille 67 -- réflexe de rédaction}

\begin{methodebox}

Avant de dériver, écrire le domaine de définition.

Avant d'appliquer Rolle, vérifier la continuité sur le segment.

Avant d'appliquer Rolle, vérifier la dérivabilité sur l'intervalle ouvert.

Avant d'appliquer le TAF, vérifier que les deux bornes appartiennent au même intervalle.

Avant d'utiliser une réciproque, vérifier la bijectivité par continuité et stricte monotonie.

Avant d'utiliser la formule de la réciproque, vérifier que la dérivée ne s'annule pas.

Pour une inégalité, construire une fonction auxiliaire dont on étudie le minimum.

Pour une unicité, chercher une stricte monotonie ou une injectivité.

Pour une existence, chercher un changement de signe ou une limite aux bornes.

Pour une borne, chercher une majoration uniforme de la dérivée.

\end{methodebox}

\subsection*{Grille 68 -- réflexe de rédaction}

\begin{methodebox}

Avant de dériver, écrire le domaine de définition.

Avant d'appliquer Rolle, vérifier la continuité sur le segment.

Avant d'appliquer Rolle, vérifier la dérivabilité sur l'intervalle ouvert.

Avant d'appliquer le TAF, vérifier que les deux bornes appartiennent au même intervalle.

Avant d'utiliser une réciproque, vérifier la bijectivité par continuité et stricte monotonie.

Avant d'utiliser la formule de la réciproque, vérifier que la dérivée ne s'annule pas.

Pour une inégalité, construire une fonction auxiliaire dont on étudie le minimum.

Pour une unicité, chercher une stricte monotonie ou une injectivité.

Pour une existence, chercher un changement de signe ou une limite aux bornes.

Pour une borne, chercher une majoration uniforme de la dérivée.

\end{methodebox}

\subsection*{Grille 69 -- réflexe de rédaction}

\begin{methodebox}

Avant de dériver, écrire le domaine de définition.

Avant d'appliquer Rolle, vérifier la continuité sur le segment.

Avant d'appliquer Rolle, vérifier la dérivabilité sur l'intervalle ouvert.

Avant d'appliquer le TAF, vérifier que les deux bornes appartiennent au même intervalle.

Avant d'utiliser une réciproque, vérifier la bijectivité par continuité et stricte monotonie.

Avant d'utiliser la formule de la réciproque, vérifier que la dérivée ne s'annule pas.

Pour une inégalité, construire une fonction auxiliaire dont on étudie le minimum.

Pour une unicité, chercher une stricte monotonie ou une injectivité.

Pour une existence, chercher un changement de signe ou une limite aux bornes.

Pour une borne, chercher une majoration uniforme de la dérivée.

\end{methodebox}

\subsection*{Grille 70 -- réflexe de rédaction}

\begin{methodebox}

Avant de dériver, écrire le domaine de définition.

Avant d'appliquer Rolle, vérifier la continuité sur le segment.

Avant d'appliquer Rolle, vérifier la dérivabilité sur l'intervalle ouvert.

Avant d'appliquer le TAF, vérifier que les deux bornes appartiennent au même intervalle.

Avant d'utiliser une réciproque, vérifier la bijectivité par continuité et stricte monotonie.

Avant d'utiliser la formule de la réciproque, vérifier que la dérivée ne s'annule pas.

Pour une inégalité, construire une fonction auxiliaire dont on étudie le minimum.

Pour une unicité, chercher une stricte monotonie ou une injectivité.

Pour une existence, chercher un changement de signe ou une limite aux bornes.

Pour une borne, chercher une majoration uniforme de la dérivée.

\end{methodebox}

\subsection*{Grille 71 -- réflexe de rédaction}

\begin{methodebox}

Avant de dériver, écrire le domaine de définition.

Avant d'appliquer Rolle, vérifier la continuité sur le segment.

Avant d'appliquer Rolle, vérifier la dérivabilité sur l'intervalle ouvert.

Avant d'appliquer le TAF, vérifier que les deux bornes appartiennent au même intervalle.

Avant d'utiliser une réciproque, vérifier la bijectivité par continuité et stricte monotonie.

Avant d'utiliser la formule de la réciproque, vérifier que la dérivée ne s'annule pas.

Pour une inégalité, construire une fonction auxiliaire dont on étudie le minimum.

Pour une unicité, chercher une stricte monotonie ou une injectivité.

Pour une existence, chercher un changement de signe ou une limite aux bornes.

Pour une borne, chercher une majoration uniforme de la dérivée.

\end{methodebox}

\subsection*{Grille 72 -- réflexe de rédaction}

\begin{methodebox}

Avant de dériver, écrire le domaine de définition.

Avant d'appliquer Rolle, vérifier la continuité sur le segment.

Avant d'appliquer Rolle, vérifier la dérivabilité sur l'intervalle ouvert.

Avant d'appliquer le TAF, vérifier que les deux bornes appartiennent au même intervalle.

Avant d'utiliser une réciproque, vérifier la bijectivité par continuité et stricte monotonie.

Avant d'utiliser la formule de la réciproque, vérifier que la dérivée ne s'annule pas.

Pour une inégalité, construire une fonction auxiliaire dont on étudie le minimum.

Pour une unicité, chercher une stricte monotonie ou une injectivité.

Pour une existence, chercher un changement de signe ou une limite aux bornes.

Pour une borne, chercher une majoration uniforme de la dérivée.

\end{methodebox}

\subsection*{Grille 73 -- réflexe de rédaction}

\begin{methodebox}

Avant de dériver, écrire le domaine de définition.

Avant d'appliquer Rolle, vérifier la continuité sur le segment.

Avant d'appliquer Rolle, vérifier la dérivabilité sur l'intervalle ouvert.

Avant d'appliquer le TAF, vérifier que les deux bornes appartiennent au même intervalle.

Avant d'utiliser une réciproque, vérifier la bijectivité par continuité et stricte monotonie.

Avant d'utiliser la formule de la réciproque, vérifier que la dérivée ne s'annule pas.

Pour une inégalité, construire une fonction auxiliaire dont on étudie le minimum.

Pour une unicité, chercher une stricte monotonie ou une injectivité.

Pour une existence, chercher un changement de signe ou une limite aux bornes.

Pour une borne, chercher une majoration uniforme de la dérivée.

\end{methodebox}

\subsection*{Grille 74 -- réflexe de rédaction}

\begin{methodebox}

Avant de dériver, écrire le domaine de définition.

Avant d'appliquer Rolle, vérifier la continuité sur le segment.

Avant d'appliquer Rolle, vérifier la dérivabilité sur l'intervalle ouvert.

Avant d'appliquer le TAF, vérifier que les deux bornes appartiennent au même intervalle.

Avant d'utiliser une réciproque, vérifier la bijectivité par continuité et stricte monotonie.

Avant d'utiliser la formule de la réciproque, vérifier que la dérivée ne s'annule pas.

Pour une inégalité, construire une fonction auxiliaire dont on étudie le minimum.

Pour une unicité, chercher une stricte monotonie ou une injectivité.

Pour une existence, chercher un changement de signe ou une limite aux bornes.

Pour une borne, chercher une majoration uniforme de la dérivée.

\end{methodebox}

\subsection*{Grille 75 -- réflexe de rédaction}

\begin{methodebox}

Avant de dériver, écrire le domaine de définition.

Avant d'appliquer Rolle, vérifier la continuité sur le segment.

Avant d'appliquer Rolle, vérifier la dérivabilité sur l'intervalle ouvert.

Avant d'appliquer le TAF, vérifier que les deux bornes appartiennent au même intervalle.

Avant d'utiliser une réciproque, vérifier la bijectivité par continuité et stricte monotonie.

Avant d'utiliser la formule de la réciproque, vérifier que la dérivée ne s'annule pas.

Pour une inégalité, construire une fonction auxiliaire dont on étudie le minimum.

Pour une unicité, chercher une stricte monotonie ou une injectivité.

Pour une existence, chercher un changement de signe ou une limite aux bornes.

Pour une borne, chercher une majoration uniforme de la dérivée.

\end{methodebox}

\subsection*{Grille 76 -- réflexe de rédaction}

\begin{methodebox}

Avant de dériver, écrire le domaine de définition.

Avant d'appliquer Rolle, vérifier la continuité sur le segment.

Avant d'appliquer Rolle, vérifier la dérivabilité sur l'intervalle ouvert.

Avant d'appliquer le TAF, vérifier que les deux bornes appartiennent au même intervalle.

Avant d'utiliser une réciproque, vérifier la bijectivité par continuité et stricte monotonie.

Avant d'utiliser la formule de la réciproque, vérifier que la dérivée ne s'annule pas.

Pour une inégalité, construire une fonction auxiliaire dont on étudie le minimum.

Pour une unicité, chercher une stricte monotonie ou une injectivité.

Pour une existence, chercher un changement de signe ou une limite aux bornes.

Pour une borne, chercher une majoration uniforme de la dérivée.

\end{methodebox}

\subsection*{Grille 77 -- réflexe de rédaction}

\begin{methodebox}

Avant de dériver, écrire le domaine de définition.

Avant d'appliquer Rolle, vérifier la continuité sur le segment.

Avant d'appliquer Rolle, vérifier la dérivabilité sur l'intervalle ouvert.

Avant d'appliquer le TAF, vérifier que les deux bornes appartiennent au même intervalle.

Avant d'utiliser une réciproque, vérifier la bijectivité par continuité et stricte monotonie.

Avant d'utiliser la formule de la réciproque, vérifier que la dérivée ne s'annule pas.

Pour une inégalité, construire une fonction auxiliaire dont on étudie le minimum.

Pour une unicité, chercher une stricte monotonie ou une injectivité.

Pour une existence, chercher un changement de signe ou une limite aux bornes.

Pour une borne, chercher une majoration uniforme de la dérivée.

\end{methodebox}

\subsection*{Grille 78 -- réflexe de rédaction}

\begin{methodebox}

Avant de dériver, écrire le domaine de définition.

Avant d'appliquer Rolle, vérifier la continuité sur le segment.

Avant d'appliquer Rolle, vérifier la dérivabilité sur l'intervalle ouvert.

Avant d'appliquer le TAF, vérifier que les deux bornes appartiennent au même intervalle.

Avant d'utiliser une réciproque, vérifier la bijectivité par continuité et stricte monotonie.

Avant d'utiliser la formule de la réciproque, vérifier que la dérivée ne s'annule pas.

Pour une inégalité, construire une fonction auxiliaire dont on étudie le minimum.

Pour une unicité, chercher une stricte monotonie ou une injectivité.

Pour une existence, chercher un changement de signe ou une limite aux bornes.

Pour une borne, chercher une majoration uniforme de la dérivée.

\end{methodebox}

\subsection*{Grille 79 -- réflexe de rédaction}

\begin{methodebox}

Avant de dériver, écrire le domaine de définition.

Avant d'appliquer Rolle, vérifier la continuité sur le segment.

Avant d'appliquer Rolle, vérifier la dérivabilité sur l'intervalle ouvert.

Avant d'appliquer le TAF, vérifier que les deux bornes appartiennent au même intervalle.

Avant d'utiliser une réciproque, vérifier la bijectivité par continuité et stricte monotonie.

Avant d'utiliser la formule de la réciproque, vérifier que la dérivée ne s'annule pas.

Pour une inégalité, construire une fonction auxiliaire dont on étudie le minimum.

Pour une unicité, chercher une stricte monotonie ou une injectivité.

Pour une existence, chercher un changement de signe ou une limite aux bornes.

Pour une borne, chercher une majoration uniforme de la dérivée.

\end{methodebox}

\subsection*{Grille 80 -- réflexe de rédaction}

\begin{methodebox}

Avant de dériver, écrire le domaine de définition.

Avant d'appliquer Rolle, vérifier la continuité sur le segment.

Avant d'appliquer Rolle, vérifier la dérivabilité sur l'intervalle ouvert.

Avant d'appliquer le TAF, vérifier que les deux bornes appartiennent au même intervalle.

Avant d'utiliser une réciproque, vérifier la bijectivité par continuité et stricte monotonie.

Avant d'utiliser la formule de la réciproque, vérifier que la dérivée ne s'annule pas.

Pour une inégalité, construire une fonction auxiliaire dont on étudie le minimum.

Pour une unicité, chercher une stricte monotonie ou une injectivité.

Pour une existence, chercher un changement de signe ou une limite aux bornes.

Pour une borne, chercher une majoration uniforme de la dérivée.

\end{methodebox}

\end{document}